\clearpage % Forces a page break and flushes out any remaining figures/tables
\section{ขอบเขตการดำเนินงาน (Scope: Features / Functions)}

\subsection{บทบาทผู้ใช้งาน (User Roles)}

\begin{tabularx}{\dimexpr\textwidth-\leftskip\relax}{@{} L{4.2cm} Y @{}}
\toprule
\hrow \thead{บทบาท (Role)} & \thead{คำอธิบายและสิทธิ์การใช้งานหลัก} \\
\midrule
นักศึกษา (Student) &
  ผู้ขอยืม: ค้นหา จองอุปกรณ์ ส่งคำขอยืม ดูสถานะคำขอ ดูประวัติ แจ้งคืน รายงานปัญหา ความเสียหาย \\
เจ้าหน้าที่ผู้ดูแล (Staff) &
  ผู้ดูแลกระบวนการ: จัดการบัญชีรายการอุปกรณ์ จัดเตรียมและส่งมอบอุปกรณ์ บันทึกการรับ–คืน ตรวจสอบสภาพ จัดการรายการเกินกำหนด \\
ผู้ดูแลระบบ (Admin) &
  ผู้ดูแลระบบโดยรวม: จัดการบัญชีผู้ใช้และสิทธิ์ ตั้งค่าระบบ และตรวจสอบบันทึกการใช้งาน \\
เจ้าหน้าที่ระดับสูง (Supervisor / อาจารย์) &
  ผู้อนุมัติระดับสูง: อนุมัติคำขอยืมอุปกรณ์ระดับ T2 (มูลค่าสูง) การต่ออายุที่เกินเงื่อนไข และการยืมของผู้ที่มีเครดิตต่ำ \\
\bottomrule
\end{tabularx}

\subsubsection*{รายละเอียดความสามารถของแต่ละบทบาท}

\noindent\textbf{ผู้ใช้ / นักศึกษา (User)} — ผู้ขอยืม (ยืมได้อย่างเดียว)
\begin{itemize}
  \item ค้นหาและดูรายการอุปกรณ์ พร้อมรายละเอียดและสถานะคงเหลือ
  \item ส่งคำขอยืม โดยระบุจำนวน เลือกหมายเลขอุปกรณ์ (T2) หรือเลือกสล็อตเวลา (T3) ตามระดับอุปกรณ์
  \item ดูเครดิตของตนเอง สิทธิ์การยืม และเงื่อนไขที่เกี่ยวข้อง
  \item ติดตามสถานะคำขอ (รออนุมัติ อนุมัติ ปฏิเสธ กำลังยืม เกินกำหนด)
  \item ถ่ายรูปตอนรับและตอนคืน อัปโหลดเข้าระบบ และกดยืนยันการรับ–คืน
  \item ขอเปลี่ยนอุปกรณ์ได้ตอนรับ (เฉพาะ T1) และขอต่ออายุการยืม (ออนไลน์ หรือ นำอุปกรณ์มาให้เจ้าหน้าที่ตรวจสภาพตามรอบ)
  \item แจ้งคืน แจ้งปัญหา และรายงานความเสียหายเบื้องต้น พร้อมดูประวัติการยืมของตนเอง
  \item \emph{ข้อจำกัด:} ยืมได้อย่างเดียว ไม่มีสิทธิ์จัดการอุปกรณ์หรืออนุมัติคำขอ
\end{itemize}

\noindent\textbf{เจ้าหน้าที่ผู้ดูแล (Staff)} — ผู้ให้ยืมและดูแลอุปกรณ์
\begin{itemize}
  \item มีสิทธิ์ยืมอุปกรณ์ได้เช่นเดียวกับผู้ใช้
  \item ลงทะเบียนและจัดการบัญชีรายการอุปกรณ์ที่เปิดให้ยืม (เพิ่ม/แก้ไขรายการ จำนวน และหมายเลขประจำอุปกรณ์)
  \item จัดเตรียมและเลือกชิ้นอุปกรณ์ที่ส่งมอบ พร้อมบันทึกหมายเลข (T1, T2)
  \item บันทึกการรับ–คืน ตรวจสอบหมายเลข และยืนยันรายการ
  \item ตรวจสอบสภาพอุปกรณ์ เทียบรูปตอนรับกับตอนคืน และประเมิน/บันทึกความเสียหาย (ระดับ B0–B3)
  \item ตรวจสภาพอุปกรณ์เมื่อผู้ยืมนำมาต่ออายุตามรอบ (การต่อออนไลน์สลับกับการตรวจสภาพ)
  \item จัดการรายการเกินกำหนดและติดตามการคืน
  \item ตรวจสอบความเรียบร้อยของสถานที่/อุปกรณ์ติดที่วันละครั้ง (T3)
\end{itemize}

\noindent\textbf{อาจารย์ (Supervisor)} — ผู้อนุมัติระดับสูง
\begin{itemize}
  \item มีสิทธิ์ยืมอุปกรณ์ได้เช่นเดียวกับผู้ใช้
  \item อนุมัติหรือปฏิเสธคำขอยืมอุปกรณ์ระดับ T2 (มูลค่าสูง)
  \item อนุมัติการต่ออายุที่เกินเงื่อนไขปรกติ (เช่น การต่ออายุครั้งถัดไปของ T1 และการต่ออายุ T2)
  \item อนุมัติการยืมของผู้ที่มีเครดิตต่ำตามที่ระบบกำหนด (เช่น D2: T1, T2 และ D3: T1, T2)
  \item ใช้ดุลยพินิจในกรณีพิเศษที่ระบบไม่อนุมัติอัตโนมัติ
\end{itemize}

\noindent\textbf{ผู้ดูแลระบบ (Admin)} — ผู้ดูแลระบบโดยรวม
\begin{itemize}
  \item จัดการบัญชีผู้ใช้และกำหนดบทบาท/สิทธิ์ โดยเฉพาะการกำหนดว่าผู้ใดเป็นเจ้าหน้าที่
  \item ตั้งค่าระบบโดยรวม เช่น พารามิเตอร์ของระบบเครดิต การจัดระดับอุปกรณ์ และเงื่อนไขสิทธิ์การยืม
  \item ตรวจสอบบันทึกการใช้งานของระบบ (Audit Log)
  \item ดูแลความปลอดภัยและการควบคุมการเข้าถึงระบบ
\end{itemize}

\noindent\textit{หมายเหตุ:} ทุกบทบาทเข้าสู่ระบบด้วยบัญชีอีเมลของมหาวิทยาลัย (KU e-mail)

\noindent\textit{หมายเหตุ (ขอบเขตอำนาจ):} เจ้าหน้าที่ (Staff) และอาจารย์ (Supervisor) มีขอบเขตอำนาจตามภาควิชาหรือหน่วยงานของตนเอง ไม่สามารถอนุมัติหรือให้ยืมอุปกรณ์ข้ามภาควิชาได้ เช่น อาจารย์ภาควิชาวิศวกรรมคอมพิวเตอร์ไม่สามารถอนุมัติการยืมอุปกรณ์ของภาควิชาวิศวกรรมไฟฟ้า

\subsection{แผนภาพยูสเคส (Use Case Diagram)}

แผนภาพยูสเคสแสดงมุมมองเชิงหน้าที่ของระบบ ว่าผู้กระทำ (Actor) แต่ละบทบาทมีปฏิสัมพันธ์กับระบบผ่านยูสเคสใดบ้าง ตารางในหัวข้อนี้สรุปยูสเคสแยกตามบทบาท ส่วนความสัมพันธ์ระหว่างยูสเคสแสดงในรูปที่~\ref{fig:usecase} โดยความสัมพันธ์มีสองแบบ ได้แก่

\begin{itemize}
  \item \emph{«include»} — ยูสเคสที่ถูกเรียกใช้เป็นส่วนหนึ่งเสมอเมื่อทำยูสเคสหลัก เช่น การส่งคำขอยืมจะตรวจสิทธิ์และเครดิตทุกครั้ง
  \item \emph{«extend»} — ยูสเคสที่ทำงานเพิ่มเติมเฉพาะในบางเงื่อนไข เช่น อาจารย์อนุมัติคำขอเฉพาะอุปกรณ์ระดับ T2 หรือการรายงานความเสียหายเมื่อพบระหว่างตรวจสอบหลังคืน
\end{itemize}

\noindent\textbf{ผู้ยืม (User)}

\begin{tabularx}{\dimexpr\textwidth-\leftskip\relax}{@{} L{4.6cm} Y @{}}
\toprule
\hrow \thead{ยูสเคส} & \thead{รายละเอียดและความสัมพันธ์} \\
\midrule
ค้นหา / กรองอุปกรณ์ & ค้นหาและกรองรายการอุปกรณ์ พร้อมดูรายละเอียดและสถานะคงเหลือ \\
ส่งคำขอยืม & สร้างคำขอยืม โดย «include» ยูสเคส \emph{ตรวจสิทธิ์ + เครดิต} ทุกครั้ง \\
ติดตามสถานะคำขอ & ดูสถานะคำขอ: รออนุมัติ อนุมัติ ปฏิเสธ กำลังยืม เกินกำหนด \\
ขอต่ออายุการยืม & ขอต่ออายุแบบออนไลน์ หรือผ่านการตรวจสภาพโดยเจ้าหน้าที่ตามรอบ \\
รับอุปกรณ์ & ยืนยันการรับ โดย «include» ยูสเคส \emph{อัปโหลดรูปหลักฐาน} (รูปตอนรับ) \\
คืนอุปกรณ์ & แจ้งคืน โดย «include» ยูสเคส \emph{อัปโหลดรูปหลักฐาน} (รูปตอนคืน) \\
ยื่นอุทธรณ์ผลประเมิน & อุทธรณ์ผลประเมินความเสียหาย โดย «include» \emph{อัปโหลดรูปหลักฐาน} (ก่อน–หลัง) \\
ดูเครดิต / ประวัติ & ดูคะแนนเครดิต ระดับเครดิต และประวัติการยืมของตนเอง \\
\bottomrule
\end{tabularx}

\vspace{0.6em}
\noindent\textbf{เจ้าหน้าที่ผู้ดูแล (Staff)}

\begin{tabularx}{\dimexpr\textwidth-\leftskip\relax}{@{} L{4.6cm} Y @{}}
\toprule
\hrow \thead{ยูสเคส} & \thead{รายละเอียดและความสัมพันธ์} \\
\midrule
จัดเตรียม / เลือก Serial & จัดเตรียมและเลือกชิ้นอุปกรณ์ที่ส่งมอบ พร้อมผูกหมายเลข (T1, T2) \\
บันทึกการส่งมอบ & บันทึกการส่งมอบและหมายเลขประจำอุปกรณ์ที่ส่งให้ผู้ยืม \\
ตรวจสภาพเพื่อต่ออายุ & ตรวจสภาพอุปกรณ์เมื่อผู้ยืมนำมาต่ออายุตามรอบ (สลับกับการต่อออนไลน์) \\
ตรวจสอบหลังคืน & ตรวจสภาพอุปกรณ์เมื่อรับคืน — «extend» ด้วย \emph{รายงานความเสียหาย} เมื่อพบความเสียหาย \\
รายงานความเสียหาย & บันทึกและประเมินระดับความเสียหาย (B0–B3) ต่อยอดจากการตรวจสอบหลังคืน \\
จัดการข้อมูลอุปกรณ์ & ลงทะเบียนและแก้ไขบัญชีรายการอุปกรณ์ จำนวน และหมายเลขประจำอุปกรณ์ \\
\bottomrule
\end{tabularx}

\vspace{0.6em}
\noindent\textbf{อาจารย์ (Supervisor)}

\begin{tabularx}{\dimexpr\textwidth-\leftskip\relax}{@{} L{4.6cm} Y @{}}
\toprule
\hrow \thead{ยูสเคส} & \thead{รายละเอียดและความสัมพันธ์} \\
\midrule
อนุมัติ / ปฏิเสธคำขอ T2 & พิจารณาอนุมัติหรือปฏิเสธคำขอ — «extend» ยูสเคส \emph{ส่งคำขอยืม} เฉพาะอุปกรณ์ระดับ T2 \\
อนุมัติต่ออายุ T2 & อนุมัติการต่ออายุที่เกินเงื่อนไข — «extend» ยูสเคส \emph{ขอต่ออายุการยืม} \\
ตัดสินอุทธรณ์ & ชี้ขาดผลอุทธรณ์ — «extend» ยูสเคส \emph{ยื่นอุทธรณ์ผลประเมิน} \\
แก้ไขผลประเมิน & ปรับระดับความเสียหายตามผลการตัดสินอุทธรณ์ \\
อนุมัติปลดระวาง & อนุมัติการปลดระวางอุปกรณ์ตามที่เจ้าหน้าที่เสนอ พร้อมบันทึก Audit Log \\
\bottomrule
\end{tabularx}

\vspace{0.6em}
\noindent\textbf{ผู้ดูแลระบบ (Admin)}

\begin{tabularx}{\dimexpr\textwidth-\leftskip\relax}{@{} L{4.6cm} Y @{}}
\toprule
\hrow \thead{ยูสเคส} & \thead{รายละเอียดและความสัมพันธ์} \\
\midrule
จัดการบัญชีผู้ใช้ & จัดการบัญชีผู้ใช้และกำหนดบทบาท/สิทธิ์ โดยเฉพาะการกำหนดเจ้าหน้าที่ \\
ตั้งค่าระบบ / เครดิต & ตั้งค่าพารามิเตอร์ของระบบ การจัดระดับอุปกรณ์ และระบบเครดิต \\
ดู Dashboard / สถิติ & ดูภาพรวมและรายงานสถิติการใช้งานของระบบ \\
\bottomrule
\end{tabularx}

% =============================================================================
%  fig_usecase.tex — figure float wrapping the UML use-case diagram.
%  The diagram body lives in _usecase_tikz.tex (shared with the standalone
%  test wrapper). Scaled to the text width; the diagram's aspect ratio leaves
%  room for the caption on a dedicated float page.
% =============================================================================
\begin{figure}[p]
  \centering
  \resizebox{\textwidth}{!}{%
    % =============================================================================
%  _usecase_tikz.tex — UML Use Case diagram (TikZ body only)
%  Recreated from resource/ULMs_Software_Project_Presentation_V.pdf (slide 1).
%  Requires (loaded in the preamble of whichever document \inputs this):
%     fontspec (Laksaman), xcolor, tikz + libraries
%       shapes.geometric, arrows.meta, positioning, calc, fit, backgrounds
%  Grayscale palette so it stays on-theme with the monochrome document; the
%  four actors are separated by layout (left = Staff/Supervisor, right =
%  User/Admin) rather than by colour.
% =============================================================================

% ---- monochrome palette (redefined here so the file is portable) ------------
\definecolor{ucFill}{HTML}{ECECEC}   % use-case ellipse fill (light gray)
\definecolor{ucEdge}{HTML}{4A4A4A}   % borders / association lines
\definecolor{ucRel} {HTML}{2E2E2E}   % «include» / «extend» dashed relations

% A simple stick-figure actor, placed via \pic at a coordinate.
\tikzset{
  actor/.pic={
    \draw[line width=0.5pt]
      (0,0.42) circle (0.17)          % head
      (0,0.25) -- (0,-0.28)           % torso
      (-0.28,0.02) -- (0.28,0.02)     % arms
      (0,-0.28) -- (-0.22,-0.62)      % left leg
      (0,-0.28) -- (0.22,-0.62);      % right leg
  }
}

\begin{tikzpicture}[
    x=1cm, y=1cm,
    every node/.style = {font=\scriptsize},
    uc/.style       = {ellipse, draw=ucEdge, line width=0.4pt, fill=ucFill,
                       align=center, text width=2.35cm, minimum width=2.6cm,
                       minimum height=0.9cm, inner sep=1.5pt},
    assoc/.style    = {draw=ucEdge!75, line width=0.3pt},
    rel/.style      = {draw=ucRel, dash pattern=on 2.2pt off 1.6pt,
                       -{Stealth[length=1.8mm]}, line width=0.4pt},
    rlbl/.style     = {font=\tiny\itshape, fill=white, inner sep=0.7pt, text=ucRel},
    boundary/.style = {draw=ucEdge, line width=0.7pt, rounded corners=4pt},
    actorlbl/.style = {font=\scriptsize, align=center, text width=2.1cm}
  ]

  % ---------------- system boundary ----------------------------------------
  \draw[boundary] (1.5,-1.95) rectangle (15.2,22.2);
  \node[font=\small\bfseries, align=center] at (8.35,21.55)
    {ระบบบริหารจัดการการยืม–คืนอุปกรณ์ (ULMs)};

  % ---------------- use cases: Staff (left column) -------------------------
  \node[uc] (s1) at (3.2,20.0) {จัดเตรียม / เลือก Serial};
  \node[uc] (s2) at (3.2,18.5) {บันทึกการส่งมอบ};
  \node[uc] (s3) at (3.2,17.0) {ตรวจสภาพเพื่อต่ออายุ};
  \node[uc] (s4) at (3.2,15.5) {จัดการข้อมูลอุปกรณ์};
  \node[uc] (s5) at (3.2,14.0) {ตรวจสอบหลังคืน};
  \node[uc] (s6) at (6.9,14.0) {รายงานความเสียหาย};

  % ---------------- use cases: Supervisor / อาจารย์ (left, lower) ----------
  \node[uc] (v1) at (3.2,10.8) {อนุมัติ / ปฏิเสธคำขอ T2};
  \node[uc] (v2) at (3.2,9.3)  {อนุมัติต่ออายุ T2};
  \node[uc] (v3) at (3.2,7.8)  {ตัดสินอุทธรณ์};
  \node[uc] (v4) at (3.2,6.3)  {แก้ไขผลประเมิน};
  \node[uc] (v5) at (3.2,4.8)  {อนุมัติปลดระวาง};

  % ---------------- use cases: User / ผู้ยืม (centre column) ---------------
  \node[uc] (u1) at (11.0,13.5) {ค้นหา / กรองอุปกรณ์};
  \node[uc] (u2) at (11.0,12.0) {ส่งคำขอยืม};
  \node[uc] (u3) at (11.0,10.5) {ติดตามสถานะคำขอ};
  \node[uc] (u4) at (11.0,9.0)  {ขอต่ออายุการยืม};
  \node[uc] (u5) at (11.0,7.5)  {รับอุปกรณ์};
  \node[uc] (u6) at (11.0,6.0)  {ยื่นอุทธรณ์ผลประเมิน};
  \node[uc] (u7) at (11.0,4.5)  {คืนอุปกรณ์};
  \node[uc] (u8) at (11.0,3.0)  {ดูเครดิต / ประวัติ};

  % ---------------- «include» targets (shared) ----------------------------
  \node[uc] (i1) at (6.9,12.6) {ตรวจสิทธิ์ + เครดิต};
  \node[uc] (i2) at (6.9,5.2)  {อัปโหลดรูปหลักฐาน};

  % ---------------- use cases: Admin / ผู้ดูแลระบบ (right, lower) ----------
  \node[uc] (a1) at (13.0,1.6)  {จัดการบัญชีผู้ใช้};
  \node[uc] (a2) at (13.0,0.1)  {ตั้งค่าระบบ / เครดิต};
  \node[uc] (a3) at (13.0,-1.4) {ดู Dashboard / สถิติ};

  % ---------------- actors --------------------------------------------------
  \coordinate (aStaff) at (-0.2,17.0);
  \pic at (aStaff) {actor};
  \node[actorlbl, below=0.55cm of aStaff] {เจ้าหน้าที่\\(Staff)};

  \coordinate (aSup) at (-0.2,7.8);
  \pic at (aSup) {actor};
  \node[actorlbl, below=0.55cm of aSup] {อาจารย์\\(Supervisor)};

  \coordinate (aUser) at (17.6,10.5);
  \pic at (aUser) {actor};
  \node[actorlbl, below=0.55cm of aUser] {ผู้ยืม\\(User)};

  \coordinate (aAdmin) at (17.6,0.1);
  \pic at (aAdmin) {actor};
  \node[actorlbl, below=0.55cm of aAdmin] {ผู้ดูแลระบบ\\(Admin)};

  % ---------------- associations: Staff ------------------------------------
  \foreach \n in {s1,s2,s3,s4,s5,s6}
    \draw[assoc] (aStaff) -- (\n.west);

  % ---------------- associations: Supervisor -------------------------------
  \foreach \n in {v1,v2,v3,v4,v5}
    \draw[assoc] (aSup) -- (\n.west);

  % ---------------- associations: User -------------------------------------
  \foreach \n in {u1,u2,u3,u4,u5,u6,u7,u8}
    \draw[assoc] (aUser) -- (\n.east);

  % ---------------- associations: Admin ------------------------------------
  \foreach \n in {a1,a2,a3}
    \draw[assoc] (aAdmin) -- (\n.east);

  % ---------------- «include» relations ------------------------------------
  % Labels placed near the source use case (pos=0.25) so the three arrows that
  % converge on อัปโหลดรูปหลักฐาน don't stack their labels on top of each other.
  \draw[rel] (u2.west) -- (i1.east)       node[rlbl, pos=0.45, above]{«include»};
  \draw[rel] (u5.west) -- (i2.north east) node[rlbl, pos=0.25, above]{«include»};
  \draw[rel] (u6.west) -- (i2.east)       node[rlbl, pos=0.25, above]{«include»};
  \draw[rel] (u7.west) -- (i2.south east) node[rlbl, pos=0.25, below]{«include»};

  % ---------------- «extend» relations -------------------------------------
  % Staff: รายงานความเสียหาย extends ตรวจสอบหลังคืน
  \draw[rel] (s6.west) -- (s5.east) node[rlbl, midway, above]{«extend»};
  % Supervisor approvals extend the User use cases they gate
  \draw[rel] (v1.east) -- (u2.west) node[rlbl, pos=0.6, above]{«extend»};
  \draw[rel] (v2.east) -- (u4.west) node[rlbl, pos=0.6, above]{«extend»};
  \draw[rel] (v3.east) -- (u6.west) node[rlbl, pos=0.6, below]{«extend»};

\end{tikzpicture}
%
  }
  \caption{แผนภาพยูสเคสของระบบ ULMs แสดงผู้กระทำ (Actor) ทั้งสี่บทบาทและยูสเคสที่เกี่ยวข้อง
    ความสัมพันธ์แบบ \emph{«include»} คือยูสเคสที่ถูกเรียกใช้เสมอ (เช่น การส่งคำขอยืมต้องตรวจสิทธิ์และเครดิตทุกครั้ง)
    และแบบ \emph{«extend»} คือยูสเคสที่ทำงานเพิ่มในบางเงื่อนไข (เช่น อาจารย์อนุมัติคำขอ T2 หรือการรายงานความเสียหายเมื่อพบระหว่างตรวจสอบหลังคืน)}
  \label{fig:usecase}
\end{figure}


\subsection{กระบวนการทำงานของระบบ (System Workflow)}

\begin{enumerate}
  \item นักศึกษาเข้าสู่ระบบ ค้นหาอุปกรณ์ และส่งคำขอยืม (ระบุรายการและช่วงเวลา)
  \item ระบบหรืออาจารย์(กรณีของระดับสูง)ประเมิณสิทธ์ในการยืมแล้วอนุมัติหรือปฏิเสธคำขอ
  \item เจ้าหน้าที่ตรวจสอบและจัดเตรียมอุปกรณ์ที่อนุมัติ พร้อมบันทึกหมายเลขประจำอุปกรณ์ (T1, T2) หรือยืนยันช่วงเวลาที่จอง (T3)
  \item เมื่ออนุมัติ นักศึกษามารับอุปกรณ์ นักเรียนกดยืนยัน เพื่อบันทึกการรับ
  \item นักศึกษาใช้งานอุปกรณ์ตามปรกติ ระบบจะแจ้งเตือนเมื่อใกล้ถึงกำหนดคืน
  \item นักศึกษาคืนอุปกรณ์ เจ้าหน้าที่ตรวจสอบสภาพและยืนยัน เพื่อบันทึกการคืน
  \item หากพบความเสียหาย เจ้าหน้าที่บันทึกรายการความเสียหาย จากนั้นปิดรายการ
\end{enumerate}

โดยกระบวนการโดยรวม รวมถึงการอนุมัติที่แตกต่างกันตามระดับอุปกรณ์ การต่ออายุ และการเชื่อมกับระบบเครดิต แสดงในรูปที่~\ref{fig:borrowflow}

% =============================================================================
%  fig_borrow_flow.tex — figure float wrapping the borrow-rule flowchart.
%  The diagram body lives in _borrow_flow_tikz.tex (shared with the standalone
%  test wrapper). Scaled to fit the text height, leaving room for the caption.
% =============================================================================
\begin{figure}[p]
  \centering
  \resizebox{!}{0.86\textheight}{%
    % =============================================================================
%  _borrow_flow_tikz.tex — the borrow-rule flowchart (TikZ body only)
%  Recreated from resource/Flowchart Whiteboard ... Neon Style-1.pdf
%  Requires (loaded in the preamble of whichever document \inputs this):
%     fontspec (Laksaman), xcolor, tikz + libraries
%       shapes.geometric, arrows.meta, positioning, calc, fit, backgrounds
%  Colours are (re)defined here so the file is portable between the proposal
%  and the standalone test wrapper.
% =============================================================================

% ---- light, restrained palette (a calmer take on the neon reference) --------
\definecolor{flStart}{HTML}{CDEFD6}   % green   — start / terminal
\definecolor{flProc} {HTML}{E4EEFB}   % blue    — borrower steps
\definecolor{flStaff}{HTML}{D3EFEA}   % teal    — staff (ผู้ให้ยืม) steps
\definecolor{flDec}  {HTML}{FBF0CC}   % yellow  — decisions
\definecolor{flUse}  {HTML}{FBDFC2}   % orange  — in-use
\definecolor{flEnd}  {HTML}{F6D3E7}   % pink    — final update
\definecolor{flStop} {HTML}{F8CFCF}   % red     — try again
\definecolor{flEdge} {HTML}{5A5A5A}   % node borders / arrows

\begin{tikzpicture}[
    node distance = 6.5mm and 10mm,
    every node/.style = {font=\footnotesize},
    box/.style       = {rectangle, rounded corners=2pt, draw=flEdge, line width=0.4pt,
                        align=center, text width=3.2cm, inner sep=3.5pt, minimum height=8mm},
    startstop/.style = {box, fill=flStart, text width=2.7cm},
    proc/.style      = {box, fill=flProc},
    staff/.style     = {box, fill=flStaff},
    inuse/.style     = {box, fill=flUse, text width=2.7cm},
    endnode/.style   = {box, fill=flEnd, text width=2.9cm},
    stopnode/.style  = {box, fill=flStop, text width=1.9cm, minimum height=7mm},
    decision/.style  = {diamond, aspect=2.1, draw=flEdge, line width=0.4pt, fill=flDec,
                        align=center, inner sep=0pt, text width=2.3cm},
    note/.style      = {rectangle, rounded corners=2pt, draw=flEdge!60, dash pattern=on 2pt off 1.5pt,
                        fill=white, align=left, text width=4.5cm, inner sep=4pt, font=\scriptsize},
    arr/.style       = {-{Stealth[length=2mm]}, draw=flEdge, line width=0.5pt},
    lbl/.style       = {font=\scriptsize, fill=white, inner sep=1pt}
  ]

  % ---------------- main spine (ผู้ยืม / Borrower) --------------------------
  \node[startstop]                     (start)  {เลือกของที่จะยืม};
  \node[proc, below=of start]          (tier)   {จำแนกระดับอุปกรณ์\\T0 \textperiodcentered\ T1 \textperiodcentered\ T2 \textperiodcentered\ T3};
  \node[proc, below=of tier]           (req)    {ส่งคำขอยืม\\(ระบุจำนวน / หมายเลข / ช่วงเวลา ตามระดับ)};
  \node[decision, below=8mm of req]    (auth)   {การอนุมัติ\\ผ่านหรือไม่?};
  \node[staff, below=8mm of auth]      (prep)   {เจ้าหน้าที่จัดเตรียม / เลือกอุปกรณ์\\(T1: ชิ้นพร้อมใช้ \textperiodcentered\ T2: ตามหมายเลข)};
  \node[proc, below=of prep]           (pin)    {ถ่ายรูปตอนรับ + อัปโหลดเข้าระบบ};
  \node[inuse, below=of pin]           (use)    {กำลังใช้งาน\\(IN USE)};
  \node[decision, below=7mm of use]    (renew)  {ต้องการ\\ต่ออายุ?};
  \node[decision, right=14mm of renew, text width=2.2cm] (ronline) {เหลือสิทธิ์\\ต่อออนไลน์?};
  \node[proc,  above=16mm of ronline, text width=2.8cm] (rsys) {ต่อผ่านระบบ (ออนไลน์) 1 ครั้ง};
  \node[staff, below=7mm of ronline, text width=2.8cm] (rchk) {นำอุปกรณ์มาให้เจ้าหน้าที่ตรวจสภาพ};
  \node[proc, below=7mm of renew]      (pout)   {ถ่ายรูปตอนคืน + อัปโหลดเข้าระบบ};
  \node[decision, below=7mm of pout]   (insp)   {ตรวจสอบหลังคืน\\พบความเสียหาย?};
  \node[endnode, below=8mm of insp]    (upd)    {อัปเดตข้อมูล\\ของที่ยืมและผู้ยืม \; (จบ)};

  % ---------------- side branches ------------------------------------------
  \node[stopnode, left=14mm of auth]   (retry)  {Try Again};
  \node[staff, right=12mm of insp, text width=3.0cm] (dmg) {รายงานความเสียหาย\\$\rightarrow$ หักเครดิต (discredit)};

  % tier-specific approval note
  \node[note, right=8mm of auth] (authnote)
    {\textbf{การอนุมัติตามระดับ}\\[1pt]
     T0–T1: ระบบตรวจสิทธิ์ + เครดิต + ของคงเหลือ\\
     T2: อาจารย์ (Supervisor) อนุมัติ\\
     T3: ตรวจช่วงเวลาว่าง (IsFree)};

  % ---------------- arrows: main spine -------------------------------------
  \draw[arr] (start) -- (tier);
  \draw[arr] (tier)  -- (req);
  \draw[arr] (req)   -- (auth);
  \draw[arr] (auth)  -- node[lbl, right]{Accept} (prep);
  \draw[arr] (prep)  -- (pin);
  \draw[arr] (pin)   -- (use);
  \draw[arr] (use)   -- (renew);
  \draw[arr] (renew) -- node[lbl, right]{ไม่} (pout);
  \draw[arr] (pout)  -- (insp);
  \draw[arr] (insp)  -- node[lbl, right]{ไม่พบ} (upd);

  % reject -> try again -> back to request
  \draw[arr] (auth) -- node[lbl, above]{Reject} (retry);
  \draw[arr] (retry.north) |- (req.west);

  % renewal cycle: 1 online renewal, then a staff condition check, alternating
  % (a passed check re-earns the next online renewal — see caption). Both
  % outcomes MERGE onto a right-hand rail and return to IN USE via a single
  % arrow, so no return line crosses the online / staff-check nodes.
  \draw[arr] (renew.east)    -- node[lbl, above]{ใช่} (ronline.west);
  \draw[arr] (ronline.north) -- node[lbl, right]{มี} (rsys.south);
  \draw[arr] (ronline.south) -- node[lbl, right]{ไม่มี} (rchk.north);
  \coordinate (rrx) at ($(rsys.east)+(0.8,0)$);
  \draw (rsys.east) -- (rsys.east -| rrx);
  \draw (rchk.east) -- (rchk.east -| rrx);
  \draw (rsys.east -| rrx) -- (rchk.east -| rrx);
  \draw[arr] (use.east -| rrx) -- (use.east);

  % damage found -> staff report -> converge to update
  \draw[arr] (insp) -- node[lbl, above]{พบ} (dmg);
  \draw[arr] (dmg.south) |- (upd.east);

  % dashed link to the approval note
  \draw[flEdge!60, dash pattern=on 2pt off 1.5pt] (auth.east) -- (authnote.west);

\end{tikzpicture}
%
  }
  \caption{ผังงานกระบวนการยืม–คืนอุปกรณ์
    รหัสสี: สีฟ้า = ขั้นตอนของผู้ยืม, สีเขียวอมฟ้า = ขั้นตอนของ
    เจ้าหน้าที่ (Staff), สีเหลือง = จุดตัดสินใจ, สีส้ม = ระหว่างใช้งาน, สีชมพู = อัปเดตข้อมูลเมื่อจบรายการ}
  \label{fig:borrowflow}
\end{figure}


\subsection{การแบ่งระดับอุปกรณ์และกฎการยืม (Equipment Tiers \& Borrowing Rules)}
\label{subsec:tiers}

\hspace*{2em}ระบบจัดประเภทอุปกรณ์ออกเป็น 4 ระดับ (Tier) ตามมูลค่าและความเสี่ยง โดยแต่ละระดับมีขั้นตอนการอนุมัติ การส่งมอบ การตรวจความเสียหาย และการต่ออายุที่แตกต่างกัน มูลค่าเครดิตประจำอุปกรณ์แต่ละระดับเป็นมูลค่าโดยประมาณไว้เพื่อใช้คำนวณบทลงโทษในระบบเครดิต

\noindent\textbf{T0 — ยืมง่าย}\quad(ราคาต่ำกว่า 100 บาท, มูลค่าเครดิต 0 ต่อชิ้น)
\begin{itemize}
  \item \textbf{ตัวอย่าง:} สายจัมเปอร์ ตัวต้านทาน
  \item \textbf{การอนุมัติ:} ส่งคำขอยืมพร้อมจำนวน ระบบตรวจสอบสิทธิ์ (จำนวนที่ยืมได้) และของคงเหลือ แล้วอนุมัติอัตโนมัติ
  \item \textbf{ยืม–คืน:} ถ่ายรูปตอนรับและตอนคืนอุปกรณ์
  \item \textbf{ความเสียหาย:} ยอมรับตามสภาพ
  \item \textbf{การต่ออายุ:} ต่อผ่านระบบได้
\end{itemize}

\noindent\textbf{T1 — ยืมได้}\quad(ราคาไม่เกิน 1,000 บาท, มูลค่าเครดิต 5 ต่อชิ้น)
\begin{itemize}
  \item \textbf{ตัวอย่าง:} Arduino มัลติมิเตอร์
  \item \textbf{การอนุมัติ:} ยืมแบบจำกัดจำนวนชิ้น ผู้ยืมส่งคำขอ ระบบแสดงเครดิต ตรวจสอบสิทธิ์และของคงเหลือ แล้วอนุมัติอัตโนมัติ จากนั้นเจ้าหน้าที่ (staff) จัดเตรียมและเลือกชิ้นที่ส่งมอบ พร้อมบันทึกหมายเลขประจำอุปกรณ์
  \item \textbf{ยืม–คืน:} ตรวจสอบหมายเลขอุปกรณ์ ผู้ยืมกดยืนยัน และถ่ายรูปตอนรับและตอนคืน (ผู้ยืมสามารถขอเปลี่ยนอุปกรณ์ได้)
  \item \textbf{ความเสียหาย:} เจ้าหน้าที่ตรวจสอบหมายเลข เทียบรูปตอนรับกับตอนคืน และรายงานความเสียหาย (ถ้ามี) ทั้งนี้ผู้ยืมสามารถอุทธรณ์ผลการประเมินได้
  \item \textbf{การต่ออายุ:} ต่อผ่านระบบ (ออนไลน์) ได้ครั้งละ 1 ครั้ง หลังจากนั้นการต่ออายุครั้งถัดไปต้องนำอุปกรณ์มาให้เจ้าหน้าที่ตรวจสภาพก่อน จึงจะใช้งานต่อและต่อออนไลน์ได้อีก 1 ครั้ง สลับเช่นนี้ไปเรื่อยๆ (กรณีผู้ยืมมีเครดิตระดับ D2 ขึ้นไปให้ยึดตามบทในตารางบทลงโทษเครดิต)
\end{itemize}

\noindent\textbf{T2 — ยืมยาก}\quad(ราคามากกว่า 1,000 บาท, มูลค่าเครดิต 10 ต่อชิ้น)
\begin{itemize}
  \item \textbf{ตัวอย่าง:} ออสซิลโลสโคป บอร์ดราคาสูง
  \item \textbf{การอนุมัติ:} ผู้ยืมเลือกหมายเลขอุปกรณ์และส่งคำขอให้อาจารย์อนุมัติ จากนั้นเจ้าหน้าที่ (staff A) จัดเตรียมอุปกรณ์
  \item \textbf{ยืม–คืน:} ตรวจสอบหมายเลขอุปกรณ์ ผู้ยืมกดยืนยัน และถ่ายรูปตอนรับและตอนคืน
  \item \textbf{ความเสียหาย:} เจ้าหน้าที่ (staff B) ตรวจสอบหมายเลข ตรวจสภาพ และรายงานความเสียหาย (ถ้ามี) ทั้งนี้ผู้ยืมสามารถอุทธรณ์ผลการประเมินได้
  \item \textbf{การต่ออายุ:} ส่งคำขอให้อาจารย์อนุมัติ
\end{itemize}

\noindent\textbf{T3 — ของติดที่}\quad(ใช้งานภายในสถานที่ ไม่คิดมูลค่าเครดิต)
\begin{itemize}
  \item \textbf{ตัวอย่าง:} ห้องปฏิบัติการ (ห้องแลป)
  \item \textbf{การอนุมัติ:} ส่งคำขอ ดูช่วงเวลาที่ว่าง เลือกสล็อตเวลา (เป็นรายชั่วโมง) ระบบตรวจสอบสิทธิ์แล้วอนุมัติอัตโนมัติ
  \item \textbf{ยืม–คืน:} ใช้งานและคืนตามเวลาที่จอง ถ่ายรูปก่อนและหลังใช้งาน
  \item \textbf{ความเสียหาย:} เจ้าหน้าที่ตรวจสอบความเรียบร้อยวันละครั้ง
  \item \textbf{การต่ออายุ:} ตรวจสอบสล็อตเวลาและต่อเวลาผ่านระบบ
\end{itemize}

\subsubsection*{ข้อกำหนดเพิ่มเติม}
\begin{itemize}
  \item \textbf{หมายเลขประจำอุปกรณ์ (Serial Number):} ติดบนอุปกรณ์ระดับ T1–T2 เพื่อระบุแต่ละชิ้น โดยแต่ละหมายเลขจะเก็บอายุการใช้งาน ประวัติการยืม และประวัติความเสียหายที่ผ่านมา
  \item \textbf{สิทธิ์การยืม (Eligibility):} เงื่อนไขการเข้าใช้งาน เช่น ชั้นปี วิชาที่เรียน สาขาวิชา และจำนวนอุปกรณ์ที่ยืมได้
  \item \textbf{ระยะเวลาการยืม:} สูงสุด 14 วัน คืนก่อนกำหนดได้ หากต้องการยืมเกิน 14 วัน ต้องขอต่ออายุตามระดับอุปกรณ์ (Tier) และระดับเครดิต
  \item \textbf{ระยะเวลารับอุปกรณ์:} ผู้ยืมต้องมารับอุปกรณ์ภายใน 1 วันหลังจากคำขอได้รับอนุมัติ มิฉะนั้นคำขอจะถูกยกเลิกโดยอัตโนมัติ
  \item \textbf{การยืมอุปกรณ์ชิ้นเดียวกัน:} ระบบจัดการแบบ atomic ตามลำดับการส่งคำขอ (ส่งก่อนได้ก่อน) หากมีคำขอที่ขัดแย้งกัน ระบบจะยกเลิก (roll back) รายการที่ส่งช้ากว่าทั้งหมด
  \item \textbf{การยืมหลายรายการพร้อมกัน:} ผู้ยืมสามารถส่งคำขอยืมอุปกรณ์หลายชิ้นได้ในคราวเดียว โดยอุปกรณ์แต่ละชิ้นจะดำเนินไปตามกระบวนการยืมของตนเองแยกจากกัน (การอนุมัติ การจัดเตรียม การส่งมอบ และการคืน เป็นอิสระต่อกัน)
  \item \textbf{การอุทธรณ์ (Appeal):} เกิดขึ้นเมื่อผู้ยืมไม่ยอมรับผลการประเมินความเสียหายของเจ้าหน้าที่ โดยยื่นอุทธรณ์ผ่านการส่งหมายเลขประจำอุปกรณ์ รูปภาพก่อนและหลังใช้งาน ระดับการประเมินที่ได้รับ และประเด็นที่ต้องการโต้แย้ง เพื่อให้อาจารย์ตรวจสอบและตัดสิน
  \item \textbf{เครดิต (Credit):} ระดับความน่าเชื่อถือของผู้ยืม ซึ่งกำหนดสิทธิ์และเงื่อนไขการยืม (รายละเอียดในหัวข้อระบบเครดิต)
\end{itemize}


\clearpage
\subsection{ระเบียบการยืม (Availability Rules)}

\hspace*{2em}นอกเหนือจากกฎการยืมตามระดับอุปกรณ์แล้ว ระบบยังกำหนดว่าอุปกรณ์แต่ละประเภทจะ\emph{พร้อมให้ยืม}เมื่อใด โดยอุปกรณ์ทุกประเภทจะแสดง ``วันที่พร้อมให้ยืม'' (available-from) ไว้ในระบบ ซึ่งคำนวณจากวันที่เร็วที่สุดที่อุปกรณ์ประเภทนั้นจะได้รับคืนจากผู้ยืมรายก่อน บวกกับจำนวนวันที่เจ้าหน้าที่ต้องใช้จัดเตรียมอุปกรณ์ (prep days) โดยรูปแบบการแสดงผลแตกต่างกันตามระดับอุปกรณ์ ทั้งนี้ระเบียบในหัวข้อนี้เป็นฉบับเบื้องต้น (beta) ซึ่งอาจปรับปรุงรายละเอียดเพิ่มเติมในระหว่างการพัฒนา

\begin{tabularx}{\dimexpr\textwidth-\leftskip\relax}{@{} L{2.2cm} L{1.5cm} Y L{3.3cm} @{}}
\toprule
\hrow \thead{ระดับ (Tier)} & \thead{วันที่พร้อม} & \thead{รายละเอียด} & \thead{ตัวอย่าง} \\
\midrule
T0 ยืมง่าย & ไม่มี &
  ผู้ยืมอาจยืมเป็นเวลานานหรือไม่นำมาคืน และเงื่อนไขเวลาค่อนข้างผ่อนคลาย จึงคาดการณ์วันที่อุปกรณ์จะพร้อมให้ยืมได้ยาก & — \\
T1 ยืมได้ & มี &
  ระบุวันที่พร้อมให้ยืมแบบแยกตามประเภท โดยรับประกันว่ามีอย่างน้อย 1 ชิ้นพร้อมในวันที่ระบุ (ไม่เจาะจงว่าเป็นชิ้นใด) &
  Arduino 1 ตัว พร้อมให้ยืมวันที่ 22/12/2026 \\
T2 ยืมยาก & มี &
  ระบุวันที่พร้อมให้ยืมแบบแยกตามชิ้นอุปกรณ์ เนื่องจากผู้ยืมต้องเลือกหมายเลขอุปกรณ์ที่ต้องการเอง &
  ออสซิลโลสโคปหมายเลข ABC-1234 พร้อมให้ยืมวันที่ 22/11/2026 \\
T3 ของติดที่ & มี &
  ระบุเวลาที่ผู้ใช้งานปัจจุบันจะหมดเวลาใช้งานเร็วที่สุด พร้อมจำนวนที่ว่างที่จะเกิดขึ้นในเวลานั้น &
  ห้อง 7603 จะมีที่ว่าง 2 ที่ ในเวลา 14:00 น. \\
\bottomrule
\end{tabularx}

\subsection{ระเบียบการจอง (Reservation Rules)}

\hspace*{2em}ผู้ยืมสามารถจองอุปกรณ์ล่วงหน้าได้ไม่เกินวันสิ้นสุดภาคเรียนปัจจุบัน โดยแต่ละระดับอุปกรณ์มีเงื่อนไขการจองต่างกัน (ระเบียบนี้เป็นฉบับเบื้องต้น (beta) เช่นเดียวกัน) ทั้งนี้ระบบกำหนดกรณีจัดการไว้ 2 กรณี ได้แก่ กรณี G1 (มีผู้ใช้งานอยู่แต่มีผู้อื่นต้องการยืมต่อ ผู้ใช้งานเดิมจะขอต่อเวลาได้หรือไม่) และกรณี G2 (การขอยืมตามปรกติที่ช่วงเวลายืมไปคร่อมช่วงเวลาที่มีการจองไว้ก่อนแล้ว)

\begin{tabularx}{\dimexpr\textwidth-\leftskip\relax}{@{} L{2.2cm} L{1.6cm} Y Y @{}}
\toprule
\hrow \thead{ระดับ (Tier)} & \thead{จองได้} & \thead{การต่อเวลาเมื่อมีผู้ใช้งานอยู่ (G1)} & \thead{การยืมที่คร่อมช่วงที่ถูกจอง (G2)} \\
\midrule
T0 ยืมง่าย & ไม่ได้ & — & — \\
T1 ยืมได้ & ได้ &
  อนุญาตให้ยืมต่อได้ภายใต้เงื่อนไขเดียวกับกรณี G2 เนื่องจากอุปกรณ์ประเภทนี้มีจำนวนมาก ระบบจะจัดสรรชิ้นที่ว่างให้หมดก่อน จึงจะให้เลือกยืมต่อจากชิ้นที่มีผู้ใช้งานอยู่ &
  ย่นระยะเวลาการยืมของผู้ยืมรายใหม่ให้สิ้นสุดก่อนวันที่จอง บวกจำนวนวันจัดเตรียมอุปกรณ์ และแจ้งผู้ยืมก่อนตัดสินใจ \\
T2 ยืมยาก & ได้ &
  การจองเป็นเพียงการเลือกวันยืม และทุกคำขอต้องผ่านการอนุมัติของอาจารย์ จึงยึดการตัดสินใจของอาจารย์เป็นหลัก เมื่ออนุมัติแล้วให้ใช้เงื่อนไขตามกรณี G2 &
  ย่นระยะเวลาการยืมให้สิ้นสุดก่อนวันที่จอง บวกจำนวนวันจัดเตรียมอุปกรณ์ ทั้งนี้การอนุมัติให้ยืมขึ้นอยู่กับดุลยพินิจของอาจารย์ \\
T3 ของติดที่ & ได้ (ภายในวัน) &
  อนุญาตให้ยืมต่อได้ เนื่องจากผู้ใช้งานหนึ่งรายสามารถยืมพร้อมกันได้สูงสุด 2 สล็อตเวลา ซึ่งเพียงพอต่อการใช้งานอยู่แล้ว &
  ยืมตามสล็อตเวลาตามปรกติ \\
\bottomrule
\end{tabularx}


\clearpage
\subsection{ระบบเครดิต (Credit System)}
\label{subsec:credit}

\hspace*{2em}ระบบเครดิตใช้ประเมินความน่าเชื่อถือของผู้ยืม โดยผู้ยืมทุกคนเริ่มต้นที่คะแนนเต็ม 100 และจะถูกหักคะแนนเมื่อทำให้อุปกรณ์เสียหายหรือคืนล่าช้า ระดับเครดิตที่เหลืออยู่จะกำหนดระยะเวลาการยืม เงื่อนไขการต่ออายุ และการอนุมัติ ทั้งนี้ประวัติความเสียหายจะถูกเชื่อมโยงกับหมายเลขประจำอุปกรณ์เพื่อใช้อ้างอิงในอนาคต

\subsubsection*{ระดับความเสียหาย (Damage Levels)}

\hspace*{2em}ความเสียหายแบ่งออกเป็น 4 ระดับ ประเมินโดยเจ้าหน้าที่ (staff) โดยแต่ละระดับมีมูลค่าเครดิตสำหรับใช้คำนวณบทลงโทษ

\begin{tabularx}{\dimexpr\textwidth-\leftskip\relax}{@{} L{5.0cm} Y >{\centering\arraybackslash}p{1.8cm} @{}}
\toprule
\hrow \thead{ระดับความเสียหาย} & \thead{คำอธิบาย} & \thead{มูลค่าเครดิต} \\
\midrule
เสียหายจากการใช้งานปรกติ (B0) &
  ความเสียหายขนาดเล็ก การชำรุดตามอายุการใช้งาน หรือการเสียหายที่เกิดขึ้นจากการใช้งานปรกติหรืออายุขัยของอุปกรณ์ & 0 \\
เกินการใช้งานปรกติ ระดับเล็กน้อย (B1) &
  ความเสียหายที่สังเกตเห็นได้ชัด ไม่ได้เกิดจากการใช้งานปรกติ สามารถซ่อมแซมได้ หรือส่งผลต่อประสิทธิภาพการทำงานเพียงเล็กน้อย & 1 \\
เกินการใช้งานปรกติ ระดับปานกลาง (B2) &
  ความเสียหายที่ไม่ได้เกิดจากการใช้งานปรกติ ซ่อมแซมยากหรือมีต้นทุนสูง หรือส่งผลต่อประสิทธิภาพของอุปกรณ์อย่างเห็นได้ชัด & 3 \\
เกินการใช้งานปรกติ ระดับรุนแรง (B3) &
  ความเสียหายที่ไม่ได้เกิดจากการใช้งานปรกติ อุปกรณ์ชำรุดโดยสมบูรณ์ ซ่อมไม่ได้ หรือไม่สามารถใช้งานตามวัตถุประสงค์ได้อีกต่อไป & 5 \\
\bottomrule
\end{tabularx}

\vspace{0.4em}
\noindent\textit{หมายเหตุ:} สำหรับอุปกรณ์ที่เสียหายระดับ B3 (ชำรุดโดยสมบูรณ์ ไม่สามารถซ่อมหรือใช้งานได้อีก) เจ้าหน้าที่สามารถส่งคำขอปลดระวางอุปกรณ์ (Decommission) ให้อาจารย์พิจารณาอนุมัติได้ (ดูรายละเอียดในหัวข้อการดำเนินงานของเจ้าหน้าที่และผู้อนุมัติ)

\subsubsection*{การคำนวณบทลงโทษ (Penalty Calculation)}
\begin{itemize}
  \item \textbf{บทลงโทษจากความเสียหาย (discredit)} $=$ มูลค่าเครดิตอุปกรณ์ $\times$ มูลค่าเครดิตความเสียหาย
  \item \textbf{อายุของบทลงโทษ} (ระยะเวลาที่เครดิตถูกหัก) $=$ discredit $\times$ 2 วัน
  \item \textbf{บทลงโทษจากการคืนล่าช้า:} หักเครดิต $=$ (มูลค่าเครดิตอุปกรณ์ $\div$ 7) $\times$ จำนวนวันที่เกินกำหนด โดยการหักเครดิตมีผลตั้งแต่วันแรกที่เกินกำหนด แต่จะเริ่มนับอายุของบทลงโทษเมื่อนำอุปกรณ์มาคืนแล้ว
  \item \textbf{กรณีอุปกรณ์สูญหาย / ไม่นำมาคืน:} หากไม่นำอุปกรณ์มาคืนภายใน 2 สัปดาห์หลังหมดกำหนดยืม จะถือว่าอุปกรณ์สูญหาย และหักเครดิตเท่ากับบทลงโทษการคืนล่าช้าตามจำนวนวัน บวกกับ (มูลค่าเครดิตอุปกรณ์ $\times$ มูลค่าเครดิตความเสียหายระดับ B3)
  \begin{itemize}
    \item หากภายหลังนำอุปกรณ์มาคืน จะลดบทลงโทษเหลือการหักเครดิตตามสูตรการคืนล่าช้า บวกกับระดับความเสียหายจริงที่ประเมินได้
    \item หากไม่สามารถนำอุปกรณ์มาคืนได้จริง ผู้ยืมต้องติดต่อเจ้าหน้าที่เพื่อรายงานและรับทราบสถานการณ์ โดยอายุของบทลงโทษจะเริ่มนับตั้งแต่วันที่รายงาน
  \end{itemize}
\end{itemize}

\subsubsection*{ระดับเครดิตและสิทธิ์การยืม (Credit Levels)}

\begin{tabularx}{\dimexpr\textwidth-\leftskip\relax}{@{} L{2.6cm} L{1.7cm} L{1.4cm} L{3.0cm} Y @{}}
\toprule
\hrow \thead{ระดับ} & \thead{ช่วงเครดิต} & \thead{ระยะยืม (วัน)} & \thead{การต่ออายุ} & \thead{เงื่อนไขการยืม} \\
\midrule
ดี (D0) & 100–80 & 14 & ต่อออนไลน์ได้ & ตามสิทธิ์ปรกติ \\
เฝ้าระวัง (D1) & 79–50 & 7 & ต่อออนไลน์ได้ & ตามสิทธิ์ปรกติ \\
เสี่ยง (D2) & 49–30 & 7 & ต่ออายุที่เจ้าหน้าที่ (ยกเว้น T3) & ต้องขอคำอนุมัติจากอาจารย์ในการยืม T1, T2 \\
เสี่ยงสูง (D3) & $<$30 & 5 & ต้องส่งคำขอยืมใหม่ & ต้องขอคำอนุมัติจากอาจารย์ในการยืม T1, T2 \\
\bottomrule
\end{tabularx}

\subsection{ข้อมูลส่วนบุคคลที่ระบบจัดเก็บ (Personal Data Collected)}

\hspace*{2em}ระบบจัดเก็บข้อมูลส่วนบุคคลของผู้ใช้เท่าที่จำเป็นต่อการให้บริการยืม–คืนอุปกรณ์ การตรวจสอบสิทธิ์การยืม และการดำเนินงานของระบบเครดิต โดยจำแนกประเภทข้อมูลที่จัดเก็บได้ดังนี้

\begin{tabularx}{\dimexpr\textwidth-\leftskip\relax}{@{} L{5.2cm} Y @{}}
\toprule
\hrow \thead{ประเภทข้อมูล} & \thead{ตัวอย่างในระบบ} \\
\midrule
ข้อมูลระบุตัวตน (Identity) &
  ชื่อ–นามสกุล รหัสนักศึกษา และอีเมลของมหาวิทยาลัย (KU e-mail) \\
ข้อมูลคุณสมบัติการยืม (Eligibility) &
  ชั้นปี สาขาวิชา และรายวิชาที่ลงทะเบียน \\
ข้อมูลการใช้งาน (Usage) &
  ประวัติการยืม–คืน คะแนนเครดิต รายการความเสียหาย และการอุทธรณ์ \\
ข้อมูลภาพ (Images) &
  รูปถ่ายอุปกรณ์ตอนรับและตอนคืน \\
ข้อมูลบันทึกระบบ (System Logs) &
  บันทึกการใช้งาน (Audit Log) และเวลาเข้าสู่ระบบ \\
\bottomrule
\end{tabularx}

\vspace{0.4em}
\noindent\textit{หมายเหตุ:} ผู้ใช้ไม่สามารถแก้ไขข้อมูลสังกัด (เช่น คณะหรือสาขาวิชา) ของตนเองได้โดยตรง หากต้องการแก้ไขต้องติดต่อเจ้าหน้าที่ (Staff) หรือผู้ดูแลระบบ (Admin) เท่านั้น เพื่อคงความถูกต้องของข้อมูลที่ใช้ตรวจสอบสิทธิ์การยืม

\clearpage
\subsection{การดำเนินงานของเจ้าหน้าที่และผู้อนุมัติ (Staff \& Supervisor Operations)}

\hspace*{2em}หัวข้อนี้อธิบายภาระงานเชิงปฏิบัติของเจ้าหน้าที่ผู้ดูแลอุปกรณ์ (Staff) และอาจารย์ผู้อนุมัติ (Supervisor) เพื่อกำหนดขอบเขตหน้าที่และผลลัพธ์ที่ระบบต้องรองรับให้ชัดเจน ทั้งนี้เจ้าหน้าที่และอาจารย์แต่ละท่านดำเนินงานได้เฉพาะภายในภาควิชาหรือหน่วยงานที่ตนรับผิดชอบเท่านั้น

\subsubsection*{งานตั้งค่าเริ่มต้น (Setup Tasks) — เจ้าหน้าที่ (Staff)}

\hspace*{2em}งานที่ดำเนินการเป็นครั้งคราวเพื่อกำหนดข้อมูลตั้งต้นของอุปกรณ์และสถานที่ในระบบ

\begin{tabularx}{\dimexpr\textwidth-\leftskip\relax}{@{} L{2.8cm} Y L{2.8cm} L{1.8cm} @{}}
\toprule
\hrow \thead{งาน} & \thead{รายละเอียด} & \thead{ผลลัพธ์ในระบบ} & \thead{ระดับ} \\
\midrule
ลงทะเบียนประเภทอุปกรณ์ &
  กรอกชื่อและราคา กำหนดระดับ (Tier) วันจัดเตรียมอุปกรณ์ สิทธิ์การเข้าใช้ (ชั้นปี/วิชา/สาขา) และจำนวนที่ยืมพร้อมกันได้ &
  Tier, credit weight, prep\_days & ทุกระดับ \\
ลงทะเบียนชิ้นอุปกรณ์ &
  เพิ่มอุปกรณ์ทีละชิ้น โดย T2 ผูกหมายเลข serial ที่ติดบนอุปกรณ์ ส่วน T1 บันทึกเป็นจำนวน &
  ของคงเหลือ และวันที่พร้อมให้ยืม & T1, T2 \\
ลงทะเบียนสถานที่ &
  กำหนดห้อง ความยาวของสล็อตเวลา และจำนวนที่นั่ง &
  ตารางสล็อตเวลาที่ว่าง & T3 \\
กำหนดเวลารับ–คืน &
  ตั้งค่าช่วงเวลารับอุปกรณ์ เช่น 08.00–17.00 น. หากคืนช้ากว่านั้นถือเป็นการคืนช้า 1 วัน &
  work\_time & ทุกระดับ \\
ปรับสถานะชั่วคราว &
  ตั้งอุปกรณ์เป็นไม่พร้อมใช้งาน (ซ่อมบำรุงหรืองดให้ยืม) &
  อุปกรณ์หลุดจาก pool ไม่ให้ยืมและจอง & ทุกระดับ \\
\bottomrule
\end{tabularx}

\clearpage
\subsubsection*{งานประจำวัน (Daily Tasks) — เจ้าหน้าที่ (Staff)}

\hspace*{2em}งานที่เกิดขึ้นตามการดำเนินการยืม–คืนของผู้ใช้ในแต่ละวัน

\begin{tabularx}{\dimexpr\textwidth-\leftskip\relax}{@{} L{2.0cm} L{2.6cm} Y L{2.6cm} L{1.5cm} @{}}
\toprule
\hrow \thead{งาน} & \thead{เกิดจาก} & \thead{การดำเนินการ} & \thead{ผลลัพธ์} & \thead{ระดับ} \\
\midrule
จัดเตรียม &
  คำขอผ่านการอนุมัติแล้ว (โดยระบบหรืออาจารย์) &
  เลือกชิ้นที่ส่งมอบ (ระบบเสนอชิ้นที่ว่างก่อน) บันทึกหมายเลข และตรวจสภาพเบื้องต้น &
  อุปกรณ์ถูกกันไว้ (reserved) และนัดวันรับ & T1, T2 \\
ส่งมอบ &
  ถึงวัน–เวลารับที่ผู้ยืมระบุ &
  ตรวจสภาพร่วมกับผู้ยืม ผู้ยืมกดยืนยัน และถ่ายรูปตอนรับ (T1 ขอเปลี่ยนชิ้นได้) &
  สถานะ IN USE & T0–T2 \\
ยกเลิกอัตโนมัติ &
  ผู้ยืมไม่มารับภายใน 1 วันหลังอนุมัติ &
  ระบบยกเลิกคำขอโดยอัตโนมัติ &
  อุปกรณ์กลับเข้า pool & T0–T2 \\
ต่ออายุแบบตรวจสภาพ &
  T1 หมดสิทธิ์ต่อออนไลน์ หรือผู้ยืมมีเครดิตระดับ D2 &
  ตรวจสภาพอุปกรณ์ที่ผู้ยืมนำมา หากผ่านจะคืนสิทธิ์ต่อออนไลน์อีก 1 ครั้ง &
  กำหนดคืนใหม่ (ไม่เกินเพดานการต่ออายุ) & T1 (T2 ส่งอาจารย์อนุมัติ) \\
รับคืนและตรวจสภาพ &
  ผู้ยืมนำอุปกรณ์มาคืนและถ่ายรูปตอนคืน &
  ตรวจหมายเลข เทียบรูปตอนรับกับตอนคืน และประเมินความเสียหายระดับ B0–B3 &
  B0 กลับเข้า pool; B1–B3 รายงานและซ่อม & T0–T2 \\
ติดตามเกินกำหนด &
  เลยกำหนดคืน (due) แล้วยังไม่คืน &
  ติดตามผู้ยืม รับรายงานกรณีคืนไม่ได้ และทำเครื่องหมายของหายเมื่อครบ 2 สัปดาห์ &
  หักเครดิต (discredit) ตามสูตร & T1, T2 \\
ตรวจสถานที่รายวัน &
  รอบตรวจวันละครั้ง &
  ตรวจความเรียบร้อยของห้อง &
  ปกติ = ปิดรอบ; พบเสียหาย = รายงาน & T3 \\
\bottomrule
\end{tabularx}

\clearpage
\subsubsection*{งานอนุมัติและการตัดสิน (Approval \& Adjudication) — อาจารย์ (Supervisor)}

\hspace*{2em}จุดตัดสินใจที่ต้องใช้ดุลยพินิจของอาจารย์ โดยผู้อนุมัติต้องเป็นคนละคนกับผู้ขอหรือผู้ประเมิน

\begin{tabularx}{\dimexpr\textwidth-\leftskip\relax}{@{} L{3.6cm} Y L{3.4cm} @{}}
\toprule
\hrow \thead{จุดตัดสินใจ} & \thead{ข้อมูลที่ใช้พิจารณา} & \thead{ข้อบังคับ} \\
\midrule
อนุมัติคำขอยืม T2 &
  เครดิตของผู้ยืม ประวัติการยืม ประวัติความเสียหายของชิ้นนั้น และวันที่พร้อมให้ยืม &
  ผู้อนุมัติต้องไม่ใช่ผู้ขอยืม \\
อนุมัติผู้มีเครดิตต่ำ (D2/D3 ยืม T1) &
  เครดิตปัจจุบัน อุปกรณ์ที่ค้างในมือ และเหตุผลประกอบ &
  ผู้อนุมัติต้องไม่ใช่ผู้ขอยืม \\
อนุมัติต่ออายุ T2 หรือกรณีเกินเงื่อนไข &
  มีผู้จองต่อหรือไม่ หากมี จะต่อได้ถึง (วันที่จอง $-$ วันที่จัดเตรียมอุปกรณ์) &
  ยึดคิวการจองเป็นหลัก \\
ตัดสินการอุทธรณ์ &
  รูปตอนรับและตอนคืนที่วางเทียบคู่กัน ระดับที่เจ้าหน้าที่ประเมิน และข้อโต้แย้งของผู้ยืม &
  ผู้ตัดสินต้องไม่ใช่ผู้ประเมิน \\
อนุมัติปลดระวางอุปกรณ์ &
  ประวัติความเสียหาย ต้นทุนซ่อม และอายุการใช้งานของชิ้นนั้น เช่น อุปกรณ์ B3 ที่ชำรุดสมบูรณ์ หรือ B1–B2 ที่จะนำไปซ่อม &
  เจ้าหน้าที่เสนอ อาจารย์อนุมัติ พร้อมบันทึก audit log \\
\bottomrule
\end{tabularx}

\subsection{ขอบเขตฟังก์ชันที่จัดทำ (Scope Functions)}


\subsubsection*{Release 1 (R01) — ฟังก์ชันหลัก}

\begin{itemize}
  \item ระบบยืนยันตัวตนและจัดการสิทธิ์ตามบทบาท (Authentication \& Role-based Access)
  \item บัญชีรายการอุปกรณ์: ค้นหา กรอง ดูรายละเอียดและสถานะคงเหลือ
  \item การส่งคำขอยืมรายการเดี่ยวและหลายรายการ พร้อมระบุช่วงเวลา
  \item รอบกระบวนการการอนุมัติ/ปฏิเสธโดยระบบหรืออาจารย์
  \item การบันทึกการรับ–คืน และรายงานความเสียหายโดยเจ้าหน้าที่ผู้ดูแลอุปกรณ์
  \item การจัดการคลังอุปกรณ์แบบหลายหน่วยต่อประเภท
  \item การติดตามสถานะคำขอ: รออนุมัติ อนุมัติ ปฏิเสธ กำลังยืม คืนแล้ว และเกินกำหนด
  \item ประวัติการยืมรายบุคคล และ Dashboard พื้นฐาน
  \item การจัดระดับอุปกรณ์เป็น 4 ระดับ (T0–T3) พร้อมเงื่อนไขการอนุมัติและการส่งมอบตามระดับ
  \item การจัดการหมายเลขประจำอุปกรณ์ (Serial Number) สำหรับอุปกรณ์ระดับ T1–T2
  \item การแสดงวันที่พร้อมให้ยืม (available-from) และการจองล่วงหน้าแบบเลือกช่วงเวลา พร้อมปฏิทิน ตามเงื่อนไขการจองของแต่ละระดับอุปกรณ์
  \item ระบบเครดิตผู้ยืม: การหักคะแนนจากความเสียหายและการคืนล่าช้า อายุของบทลงโทษ และการปรับสิทธิ์การยืมตามระดับเครดิต
\end{itemize}

\subsubsection*{Release 2 (R02) — ฟังก์ชันเสริม}

\begin{itemize}
  \item การจัดการรายการเกินกำหนด
  \item การติดตามการซ่อมบำรุง และทำให้สถานะอุปกรณ์ไม่พร้อมใช้งานชั่วคราว
  \item การออกรายงานและส่งออกข้อมูล
  \item ระบบแจ้งเตือนและติดตามผ่าน Email รวมถึงการแจ้งเตือนภายในแอป
  \item การอุทธรณ์ผลการประเมินความเสียหายของเจ้าหน้าที่โดยอาจารย์
  \item การปลดระวางอุปกรณ์ (Decommission) โดยเจ้าหน้าที่เสนอและอาจารย์อนุมัติ พร้อมบันทึกลงใน Audit Log
\end{itemize}

\subsection{สถาปัตยกรรมระบบ (System Architecture)}


\begin{tabularx}{\dimexpr\textwidth-\leftskip\relax}{@{} L{4.2cm} Y L{3.4cm} @{}}
\toprule
\hrow \thead{ชั้น (Tier)} & \thead{หน้าที่} & \thead{เทคโนโลยี} \\
\midrule
Client (Presentation) &
  สื่อสารกับผู้ใช้งานผ่าน web application เพื่อรับ request หรือข้อมูลจากผู้ใช้งานและส่งต่อให้ API Server และในทางกลับกัน &
  React + TypeScript \\
Application Server (Business Logic) &
  ควบคุมและประมวลผลตรรกะทางธุรกิจของระบบ เช่น การยืนยันตัวตน การจัดการสิทธิ์ผู้ใช้ การตรวจสอบเงื่อนไขการยืม–คืน การอนุมัติคำขอ การแจ้งเตือน และประสานงานกับฐานข้อมูล &
  NestJS \\
API Server (API Layer) &
  เป็นช่องทางการสื่อสารระหว่าง Client กับ Application Server จัดการการรับ–ส่งข้อมูล กำหนด API Endpoint/Procedure ตรวจสอบชนิดข้อมูล (Type-safe) และส่งต่อคำขอไปยัง Business Logic &
  tRPC \\
Database Server (Data Layer) &
  จัดเก็บและบริหารข้อมูลของระบบในแบบทั้ง SQL และ NoSQL เช่น ผู้ใช้ อุปกรณ์ คำขอยืม ประวัติการยืม การแจ้งเตือน และข้อมูลการซ่อมบำรุง พร้อมรองรับการค้นหาและการทำธุรกรรม &
  PostgreSQL, MongoDB \\
\bottomrule
\end{tabularx}

\vspace{0.4em}
\noindent ลำดับการไหลของข้อมูล (Data Flow): Client $\rightarrow$ API Server $\rightarrow$ Application Server $\rightarrow$ Database Server

\clearpage
\subsection{เทคโนโลยีที่ใช้ (Technology Stack)}

\begin{tabularx}{\dimexpr\textwidth-\leftskip\relax}{@{} L{4.5cm} Y @{}}
\toprule
\hrow \thead{ด้าน} & \thead{เทคโนโลยี / เครื่องมือ} \\
\midrule
Front-End &
  \textbf{[Core]} Vite + React + TypeScript + React Router \newline
  \textbf{[UI Layer]} Tailwind CSS + shadcn/ui \newline
  \textbf{[Data/Form]} TanStack Query + React Hook Form + Zod \newline
  \textbf{[Global]} React Context \newline
  \textbf{[Scale]} Server-side Pagination/Search + Code-splitting \newline
  \textbf{[Real-time]} TanStack Query Polling \\
Back-End / API &
  \textbf{[Core]} NestJS \newline
  \textbf{[API]} tRPC \\
Database             & PostgreSQL, MongoDB, Prisma \\
Version Control      & Git + GitHub \\
Design               & Figma \\
Collaboration / PM   & Miro (Whiteboard) \\
\bottomrule
\end{tabularx}

\clearpage
\subsection{การออกแบบฐานข้อมูล (Database Design)}

ระบบจัดเก็บข้อมูลด้วยฐานข้อมูลเชิงสัมพันธ์ (Relational Database) เป็นหลัก เพื่อให้การยืม การจอง และการหักคะแนนเครดิตดำเนินอยู่ภายใต้ธุรกรรม (Transaction) เดียวกันได้ การออกแบบรอบล่าสุดปรับโครงสร้างใหม่ทั้งหมดจากรอบก่อนหน้า โดยมีแนวคิดหลักสามข้อ

\begin{itemize}
  \item \textbf{รวมอุปกรณ์และห้องเป็นทรัพยากรชนิดเดียวกัน} ผ่านตาราง \texttt{ResourceInfo} ซึ่งเป็นศูนย์กลางของระบบ ตาราง \texttt{ItemIndiv} (อุปกรณ์รายชิ้น) และ \texttt{RoomInfo} (ห้องและสถานที่) ต่างผูกกลับมาที่ทรัพยากรเดียวกัน ทำให้การยืมอุปกรณ์และการจองสถานที่เดินอยู่บนกระบวนการเดียวกันทั้งหมด
  \item \textbf{แยกกฎออกจากข้อมูล} เงื่อนไขการยืมและบทลงโทษไม่ได้ฝังอยู่ในตารางอุปกรณ์ แต่ถูกแยกเป็นชุดตารางกฎ (\texttt{BorrowRule}, \texttt{BorrowConstraints}, \texttt{CreditTier}, \texttt{PenaltyRule}) ที่ปรับค่าได้โดยไม่ต้องแก้โครงสร้าง จึงเป็นกลไกที่ใช้บังคับใช้การแบ่งระดับอุปกรณ์ในหัวข้อ \ref{subsec:tiers} และระบบเครดิตในหัวข้อ \ref{subsec:credit}
  \item \textbf{เก็บทุกเหตุการณ์เป็นบันทึก (Log)} สภาพของทรัพยากร การส่งมอบและรับคืน การซ่อม และการหักคะแนน ล้วนถูกบันทึกแยกรายการพร้อมเวลาและผู้กระทำ ทำให้ตรวจสอบย้อนหลังและใช้ประกอบการอุทธรณ์ได้
\end{itemize}

\hspace*{2em}รูปที่ \ref{fig:erd} แสดงแผนภาพความสัมพันธ์ของข้อมูล (Entity–Relationship Diagram) ของทั้งระบบ และหัวข้อย่อยถัดจากนี้คือพจนานุกรมข้อมูล (Data Dictionary) ของทั้ง 33 ตาราง โดยคอลัมน์ ``คีย์'' ระบุ PK (Primary Key) หรือ FK (Foreign Key) และคีย์หลักของทุกตารางเป็น Surrogate Key ชนิด INT ที่ระบบสร้างขึ้นเอง

% =============================================================================
%  fig_erd.tex — figure float for the Entity–Relationship diagram (§5).
%  The image is the ERD exported from dbdiagram.io, extracted from
%  resource/DatabaseDesign4.pdf (page 1) and trimmed of its white border.
%
%  Uses `sidewaysfigure` (package rotating): the diagram is landscape
%  (aspect ≈ 1.51), so rotating it onto a dedicated page lets it be scaled to
%  the FULL page height instead of the text width — roughly 1.5× larger, which
%  is what makes the ~4 pt column labels legible in print.
% =============================================================================
\begin{sidewaysfigure}
  \centering
  % Height is the binding constraint once rotated (it becomes the page width),
  % so cap both: whichever limit is hit first keeps the aspect ratio intact.
  \includegraphics[width=0.95\textheight, height=0.92\textwidth, keepaspectratio]{er_diagram.png}
  \caption{แผนภาพความสัมพันธ์ของข้อมูล (Entity–Relationship Diagram) ของระบบ ULMs
    แสดงตารางทั้งหมด 33 ตารางพร้อมความสัมพันธ์ระหว่างกัน โดยมีตาราง \texttt{ResourceInfo}
    เป็นศูนย์กลางที่เชื่อมอุปกรณ์ (\texttt{ItemIndiv}) และห้อง (\texttt{RoomInfo})
    เข้ากับการจอง การใช้งาน การตรวจสอบ และการซ่อมบำรุง}
  \label{fig:erd}
\end{sidewaysfigure}


\subsubsection*{กลุ่มบัญชีผู้ใช้และสิทธิ์ (Accounts)}

\begin{keeptable}
\noindent\textbf{AccountInfo} — ข้อมูลบัญชีผู้ใช้

\begin{tabularx}{\linewidth}{@{} L{4.3cm} L{2.1cm} L{1.0cm} Y @{}}
\toprule
\hrow \thead{ฟิลด์} & \thead{ชนิดข้อมูล} & \thead{คีย์} & \thead{คำอธิบาย} \\
\midrule
AccountKey     & INT     & PK & Surrogate Key ของตาราง \\
Email          & VARCHAR &    & อีเมลของผู้ใช้ \\
HashedPassword & VARCHAR &    & รหัสผ่านของผู้ใช้ ซึ่งต้องจัดเก็บในรูปแบบที่ผ่านการ Hash แล้วเท่านั้น \\
UserID         & VARCHAR &    & รหัสประจำตัวนักศึกษาหรือบุคลากร \\
UserFName      & VARCHAR &    & ชื่อจริงของผู้ใช้ \\
UserLName      & VARCHAR &    & นามสกุลของผู้ใช้ \\
UserCredit     & INT     &    & คะแนนเครดิตคงเหลือของบัญชี \\
RoleKey        & INT     & FK & บทบาทของผู้ใช้ในระบบ อ้างอิง \texttt{RoleInfo} \\
\bottomrule
\end{tabularx}
\end{keeptable}

\begin{keeptable}
\noindent\textbf{RoleInfo} — ตำแหน่งของผู้ใช้ในระบบ

\begin{tabularx}{\linewidth}{@{} L{4.3cm} L{2.1cm} L{1.0cm} Y @{}}
\toprule
\hrow \thead{ฟิลด์} & \thead{ชนิดข้อมูล} & \thead{คีย์} & \thead{คำอธิบาย} \\
\midrule
RoleKey  & INT     & PK & Surrogate Key ของตาราง \\
RoleName & VARCHAR &    & ชื่อบทบาทของผู้ใช้ ได้แก่ นักศึกษา (Student) / อาจารย์ (Teacher) / เจ้าหน้าที่ (Staff) / ผู้ดูแลระบบ (Admin) \\
\bottomrule
\end{tabularx}
\end{keeptable}

\begin{keeptable}
\noindent\textbf{FacultyInfo} — ข้อมูลคณะ แยกตามสาขา

\begin{tabularx}{\linewidth}{@{} L{4.3cm} L{2.1cm} L{1.0cm} Y @{}}
\toprule
\hrow \thead{ฟิลด์} & \thead{ชนิดข้อมูล} & \thead{คีย์} & \thead{คำอธิบาย} \\
\midrule
FacultyKey      & INT     & PK & Surrogate Key ของตาราง \\
FacultyName     & VARCHAR &    & ชื่อของคณะ \\
FacultyBranch   & VARCHAR &    & สาขาภายในคณะ \\
ManagementGroup & INT     & FK & กลุ่มการจัดการที่คณะนี้สังกัด อ้างอิง \texttt{ManagementGroup} \\
\bottomrule
\end{tabularx}
\end{keeptable}

\begin{keeptable}
\noindent\textbf{ManagementGroup} — กลุ่มผู้มีอำนาจให้ยืม ใช้เป็นตัวกลางเชื่อมคณะและชมรมเข้าด้วยกัน ทำให้ทั้งสองแบบใช้กลไกสิทธิ์ชุดเดียวกัน

\begin{tabularx}{\linewidth}{@{} L{4.3cm} L{2.1cm} L{1.0cm} Y @{}}
\toprule
\hrow \thead{ฟิลด์} & \thead{ชนิดข้อมูล} & \thead{คีย์} & \thead{คำอธิบาย} \\
\midrule
ManageGroupKey & INT & PK & Surrogate Key ของตาราง \\
GroupType      & INT & FK & ประเภทของกลุ่ม อ้างอิง \texttt{GroupType} \\
\bottomrule
\end{tabularx}
\end{keeptable}

\begin{keeptable}
\noindent\textbf{GroupType} — ประเภทของกลุ่มผู้ให้ยืม

\begin{tabularx}{\linewidth}{@{} L{4.3cm} L{2.1cm} L{1.0cm} Y @{}}
\toprule
\hrow \thead{ฟิลด์} & \thead{ชนิดข้อมูล} & \thead{คีย์} & \thead{คำอธิบาย} \\
\midrule
GroupTypeKey  & INT     & PK & Surrogate Key ของตาราง \\
GroupTypeName & VARCHAR &    & ชื่อประเภทของกลุ่ม ได้แก่ คณะ (Faculty) หรือ ชมรม (Club) \\
\bottomrule
\end{tabularx}
\end{keeptable}

\begin{keeptable}
\noindent\textbf{ClubInfo} — ข้อมูลชมรม

\begin{tabularx}{\linewidth}{@{} L{4.3cm} L{2.1cm} L{1.0cm} Y @{}}
\toprule
\hrow \thead{ฟิลด์} & \thead{ชนิดข้อมูล} & \thead{คีย์} & \thead{คำอธิบาย} \\
\midrule
ClubKey            & INT     & PK & Surrogate Key ของตาราง \\
ClubName           & VARCHAR &    & ชื่อของชมรม \\
ManagementGroupKey & INT     & FK & กลุ่มการจัดการที่ชมรมนี้สังกัด \\
\bottomrule
\end{tabularx}
\end{keeptable}

\begin{keeptable}
\noindent\textbf{Authority} — อำนาจของผู้ใช้ในแต่ละคณะและชมรม ผู้ใช้คนเดียวสามารถมีสิทธิ์ต่างระดับกันในแต่ละกลุ่มได้

\begin{tabularx}{\linewidth}{@{} L{4.3cm} L{2.1cm} L{1.0cm} Y @{}}
\toprule
\hrow \thead{ฟิลด์} & \thead{ชนิดข้อมูล} & \thead{คีย์} & \thead{คำอธิบาย} \\
\midrule
AuthorityKey       & INT & PK & Surrogate Key ของตาราง \\
AccountKey         & INT & FK & บัญชีผู้ใช้ที่ได้รับสิทธิ์ \\
ManagementGroupKey & INT & FK & กลุ่มที่สิทธิ์นี้มีผลบังคับ \\
AuthorityRoleKey   & INT & FK & ระดับสิทธิ์ที่ได้รับในกลุ่มนั้น \\
\bottomrule
\end{tabularx}
\end{keeptable}

\begin{keeptable}
\noindent\textbf{AuthorityRole} — ชื่อและระดับของสิทธิ์

\begin{tabularx}{\linewidth}{@{} L{4.3cm} L{2.1cm} L{1.0cm} Y @{}}
\toprule
\hrow \thead{ฟิลด์} & \thead{ชนิดข้อมูล} & \thead{คีย์} & \thead{คำอธิบาย} \\
\midrule
AuthorityRoleKey & INT     & PK & Surrogate Key ของตาราง \\
AuthorityName    & VARCHAR &    & ชื่อของสิทธิ์ ได้แก่ สมาชิก (Member) ซึ่งครอบคลุมนักศึกษาและอาจารย์ / เจ้าหน้าที่ (Staff) / ผู้อนุมัติ (Supervisor) \\
AuthorityLevel   & INT     &    & ระดับของสิทธิ์ ใช้เทียบกับ \texttt{MinimumAuthorityLevel} ในการตัดสินว่าอนุมัติได้หรือไม่ ผู้ที่ไม่มีอำนาจอนุมัติให้เป็น NULL หรือ 0 เช่น Staff = 1 และ Supervisor = 2 \\
\bottomrule
\end{tabularx}
\end{keeptable}

\subsubsection*{กลุ่มอุปกรณ์และห้อง (Item \& Room)}

\begin{keeptable}
\noindent\textbf{ItemInfo} — ข้อมูลของอุปกรณ์แต่ละประเภท

\begin{tabularx}{\linewidth}{@{} L{4.3cm} L{2.1cm} L{1.0cm} Y @{}}
\toprule
\hrow \thead{ฟิลด์} & \thead{ชนิดข้อมูล} & \thead{คีย์} & \thead{คำอธิบาย} \\
\midrule
ItemKey      & INT     & PK & Surrogate Key ของตาราง \\
ItemName     & VARCHAR &    & ชื่อของอุปกรณ์ \\
ItemDesc     & TEXT    &    & รายละเอียดของอุปกรณ์ กรณีเป็นอุปกรณ์ชิ้นเล็กที่นับรวมเป็นชุด เช่น สายไฟ 20 เส้น ให้ระบุจำนวนไว้ในฟิลด์นี้ \\
CreditWeight & FLOAT   &    & ตัวคูณน้ำหนักเครดิตของอุปกรณ์ ใช้คำนวณคะแนนที่หักเมื่อเกิดความเสียหาย \\
\bottomrule
\end{tabularx}
\end{keeptable}

\begin{keeptable}
\noindent\textbf{ItemIndiv} — ข้อมูลอุปกรณ์แยกรายชิ้น

\begin{tabularx}{\linewidth}{@{} L{4.3cm} L{2.1cm} L{1.0cm} Y @{}}
\toprule
\hrow \thead{ฟิลด์} & \thead{ชนิดข้อมูล} & \thead{คีย์} & \thead{คำอธิบาย} \\
\midrule
IndivKey    & INT     & PK & Surrogate Key ของตาราง \\
ResourceKey & INT     & FK & ทรัพยากรที่อุปกรณ์ชิ้นนี้ผูกอยู่ \\
ItemKey     & INT     & FK & อุปกรณ์ชิ้นนี้เป็นอุปกรณ์ประเภทใด \\
ItemID      & VARCHAR &    & รหัสประจำอุปกรณ์ชิ้นนั้น (หมายเลขประจำชิ้น) \\
\bottomrule
\end{tabularx}
\end{keeptable}

\begin{keeptable}
\noindent\textbf{RoomInfo} — ข้อมูลห้องและสถานที่

\begin{tabularx}{\linewidth}{@{} L{4.3cm} L{2.1cm} L{1.0cm} Y @{}}
\toprule
\hrow \thead{ฟิลด์} & \thead{ชนิดข้อมูล} & \thead{คีย์} & \thead{คำอธิบาย} \\
\midrule
RoomKey      & INT     & PK & Surrogate Key ของตาราง \\
ResourceKey  & INT     & FK & ทรัพยากรที่ห้องนี้ผูกอยู่ \\
RoomName     & VARCHAR &    & ชื่อของห้อง \\
RoomDesc     & TEXT    &    & รายละเอียดเกี่ยวกับห้อง \\
RoomLocation & VARCHAR &    & ตำแหน่งที่ตั้งของห้อง \\
CreditWeight & FLOAT   &    & ตัวคูณน้ำหนักเครดิตของห้อง \\
\bottomrule
\end{tabularx}
\end{keeptable}

\subsubsection*{กลุ่มทรัพยากรและสภาพ (Resource \& Condition)}

\begin{keeptable}
\noindent\textbf{ResourceInfo} — ตารางกลางที่รวมข้อมูลของห้องและอุปกรณ์เข้าด้วยกัน เป็นศูนย์กลางที่ตารางการยืม การจอง และการซ่อมอ้างอิงถึง

\begin{tabularx}{\linewidth}{@{} L{4.3cm} L{2.1cm} L{1.0cm} Y @{}}
\toprule
\hrow \thead{ฟิลด์} & \thead{ชนิดข้อมูล} & \thead{คีย์} & \thead{คำอธิบาย} \\
\midrule
ResourceKey   & INT  & PK & Surrogate Key ของตาราง \\
ManagedBy     & INT  & FK & กลุ่มที่เป็นผู้ดูแลทรัพยากรนี้ อ้างอิง \texttt{ManagementGroup} \\
BorrowRule    & INT  & FK & กฎการยืมที่ใช้กับทรัพยากรนี้ \\
ConditionKey  & INT  & FK & สภาพล่าสุดของทรัพยากร อ้างอิง \texttt{ConditionLog} \\
CurrentStatus & INT  & FK & สถานะปัจจุบันของทรัพยากร อ้างอิง \texttt{ResourceStatus} \\
BufferTime    & INT  &    & ระยะเวลาพักทรัพยากรระหว่างรอบการยืม หน่วยเป็นชั่วโมง \\
AllowBorrow   & BOOL &    & เปิดรับให้ยืมหรือจองเพิ่มหรือไม่ \\
\bottomrule
\end{tabularx}
\end{keeptable}

\begin{keeptable}
\noindent\textbf{ResourceStatus} — สถานะปัจจุบันของทรัพยากร

\begin{tabularx}{\linewidth}{@{} L{4.3cm} L{2.1cm} L{1.0cm} Y @{}}
\toprule
\hrow \thead{ฟิลด์} & \thead{ชนิดข้อมูล} & \thead{คีย์} & \thead{คำอธิบาย} \\
\midrule
StatusKey  & INT     & PK & Surrogate Key ของตาราง \\
StatusName & VARCHAR &    & ชื่อสถานะ เช่น อยู่ในคลัง / เตรียมของแล้ว / มอบให้ผู้ใช้แล้ว / สูญหาย \\
\bottomrule
\end{tabularx}
\end{keeptable}

\begin{keeptable}
\noindent\textbf{ConditionLog} — ประวัติการบันทึกสภาพของทรัพยากร ทุกครั้งที่มีการตรวจสอบจะเพิ่มระเบียนใหม่แทนการแก้ระเบียนเดิม

\begin{tabularx}{\linewidth}{@{} L{4.3cm} L{2.1cm} L{1.0cm} Y @{}}
\toprule
\hrow \thead{ฟิลด์} & \thead{ชนิดข้อมูล} & \thead{คีย์} & \thead{คำอธิบาย} \\
\midrule
ConditionKey & INT      & PK & Surrogate Key ของตาราง \\
ResourceKey  & INT      & FK & ทรัพยากรที่ถูกตรวจสอบ \\
LoggedBy     & INT      & FK & บัญชีผู้ใช้ที่เป็นผู้ตรวจสอบ \\
Condition    & INT      & FK & สภาพของทรัพยากรที่บันทึกไว้ อ้างอิง \texttt{ConditionType} \\
Notes        & TEXT     &    & บันทึกเพิ่มเติม \\
LoggedAt     & DATETIME &    & เวลาที่ทำการตรวจสอบ \\
\bottomrule
\end{tabularx}
\end{keeptable}

\begin{keeptable}
\noindent\textbf{ConditionType} — ประเภทสภาพของทรัพยากร

\begin{tabularx}{\linewidth}{@{} L{4.3cm} L{2.1cm} L{1.0cm} Y @{}}
\toprule
\hrow \thead{ฟิลด์} & \thead{ชนิดข้อมูล} & \thead{คีย์} & \thead{คำอธิบาย} \\
\midrule
ConditionTypeKey & INT     & PK & Surrogate Key ของตาราง \\
ConditionName    & VARCHAR &    & ชื่อสภาพของทรัพยากร เช่น ปกติ / มีรอย / ชำรุด \\
\bottomrule
\end{tabularx}
\end{keeptable}

\subsubsection*{กลุ่มกฎการยืมและบทลงโทษ (Borrow \& Penalty Rules)}

\begin{keeptable}
\noindent\textbf{BorrowRule} — กฎในการยืมทรัพยากร ทำหน้าที่เป็นหัวข้อรวมของข้อกำหนดและบทลงโทษที่ผูกอยู่ใต้กฎเดียวกัน

\begin{tabularx}{\linewidth}{@{} L{4.3cm} L{2.1cm} L{1.0cm} Y @{}}
\toprule
\hrow \thead{ฟิลด์} & \thead{ชนิดข้อมูล} & \thead{คีย์} & \thead{คำอธิบาย} \\
\midrule
BorrowRuleKey  & INT     & PK & Surrogate Key ของตาราง \\
BorrowRuleName & VARCHAR &    & ชื่อของกฎ \\
\bottomrule
\end{tabularx}
\end{keeptable}

\begin{keeptable}
\noindent\textbf{BorrowConstraints} — ข้อกำหนดของการยืม แยกตามช่วงคะแนนเครดิตของผู้ยืม กฎหนึ่งข้อจึงให้เงื่อนไขต่างกันได้ตามเครดิตของผู้ยืมแต่ละคน

\begin{tabularx}{\linewidth}{@{} L{4.3cm} L{2.1cm} L{1.0cm} Y @{}}
\toprule
\hrow \thead{ฟิลด์} & \thead{ชนิดข้อมูล} & \thead{คีย์} & \thead{คำอธิบาย} \\
\midrule
ConstraintsKey        & INT & PK & Surrogate Key ของตาราง \\
BorrowRuleKey         & INT & FK & ข้อกำหนดนี้ใช้กับกฎการยืมข้อใด \\
CreditTierKey         & INT & FK & ใช้กับผู้ยืมที่อยู่ในช่วงคะแนนเครดิตใด \\
MinimumAuthorityLevel & INT &    & ระดับสิทธิ์ขั้นต่ำของผู้ที่อนุมัติคำขอได้ หากเป็น 0 หมายถึงระบบอนุมัติได้เองโดยอัตโนมัติ \\
MaxBorrowDate         & INT &    & จำนวนวันที่ยืมได้ต่อการยืมหนึ่งครั้ง \\
MaxExtendTime         & INT &    & จำนวนครั้งที่ต่ออายุการยืมได้ ก่อนต้องนำทรัพยากรมาให้เจ้าหน้าที่ตรวจสอบ \\
\bottomrule
\end{tabularx}
\end{keeptable}

\begin{keeptable}
\noindent\textbf{CreditTier} — การแบ่งช่วงคะแนนเครดิต

\begin{tabularx}{\linewidth}{@{} L{4.3cm} L{2.1cm} L{1.0cm} Y @{}}
\toprule
\hrow \thead{ฟิลด์} & \thead{ชนิดข้อมูล} & \thead{คีย์} & \thead{คำอธิบาย} \\
\midrule
CreditTierKey  & INT     & PK & Surrogate Key ของตาราง \\
CreditTierName & VARCHAR &    & ชื่อของช่วงเครดิต \\
CreditMin      & INT     &    & คะแนนต่ำสุดของช่วง \\
CreditMax      & INT     &    & คะแนนสูงสุดของช่วง \\
\bottomrule
\end{tabularx}
\end{keeptable}

\begin{keeptable}
\noindent\textbf{PenaltyRule} — กฎในการหักคะแนนเครดิต

\begin{tabularx}{\linewidth}{@{} L{4.3cm} L{2.1cm} L{1.0cm} Y @{}}
\toprule
\hrow \thead{ฟิลด์} & \thead{ชนิดข้อมูล} & \thead{คีย์} & \thead{คำอธิบาย} \\
\midrule
PenaltyRuleKey & INT & PK & Surrogate Key ของตาราง \\
BorrowRuleKey  & INT & FK & กฎการยืมที่บทลงโทษนี้ผูกอยู่ \\
PenaltyType    & INT & FK & ประเภทของสาเหตุการลงโทษ อ้างอิง \texttt{PenaltyType} \\
PenaltyAmount  & INT &    & จำนวนคะแนนที่หัก \\
PenaltyLength  & INT &    & ระยะเวลาที่การหักคะแนนมีผลบังคับ \\
\bottomrule
\end{tabularx}
\end{keeptable}

\begin{keeptable}
\noindent\textbf{PenaltyType} — ประเภทของสาเหตุการลงโทษ

\begin{tabularx}{\linewidth}{@{} L{4.3cm} L{2.1cm} L{1.0cm} Y @{}}
\toprule
\hrow \thead{ฟิลด์} & \thead{ชนิดข้อมูล} & \thead{คีย์} & \thead{คำอธิบาย} \\
\midrule
PenaltyTypeKey & INT     & PK & Surrogate Key ของตาราง \\
PenaltyType    & VARCHAR &    & ชื่อประเภทของสาเหตุ เช่น ทรัพยากรชำรุด / ทรัพยากรสูญหาย \\
\bottomrule
\end{tabularx}
\end{keeptable}

\begin{keeptable}
\noindent\textbf{Eligibility} — กฎว่าผู้ใช้จากกลุ่มใดและตำแหน่งใดมีสิทธิ์ยืมทรัพยากรใดได้บ้าง เป็นกลไกที่ใช้บังคับขอบเขตการให้ยืมข้ามคณะ

\begin{tabularx}{\linewidth}{@{} L{4.3cm} L{2.1cm} L{1.0cm} Y @{}}
\toprule
\hrow \thead{ฟิลด์} & \thead{ชนิดข้อมูล} & \thead{คีย์} & \thead{คำอธิบาย} \\
\midrule
EliKey      & INT & PK & Surrogate Key ของตาราง \\
ResourceKey & INT & FK & ทรัพยากรที่กำหนดสิทธิ์ \\
GroupKey    & INT & FK & กลุ่มที่มีสิทธิ์ยืม อ้างอิง \texttt{ManagementGroup} \\
RoleKey     & INT & FK & ตำแหน่งภายในกลุ่มที่มีสิทธิ์ยืม อ้างอิง \texttt{AuthorityRole} \\
\bottomrule
\end{tabularx}
\end{keeptable}

\subsubsection*{กลุ่มการจอง (Reservation)}

\begin{keeptable}
\noindent\textbf{Reservations} — ตารางเวลาการจองทรัพยากรล่วงหน้า

\begin{tabularx}{\linewidth}{@{} L{4.3cm} L{2.1cm} L{1.0cm} Y @{}}
\toprule
\hrow \thead{ฟิลด์} & \thead{ชนิดข้อมูล} & \thead{คีย์} & \thead{คำอธิบาย} \\
\midrule
ReservationKey    & INT      & PK & Surrogate Key ของตาราง \\
ResourceKey       & INT      & FK & ทรัพยากรที่ถูกจอง \\
ReservedBy        & INT      & FK & บัญชีผู้ใช้ที่เป็นผู้จอง \\
Reason            & VARCHAR  &    & เหตุผลโดยย่อของการจอง เช่น ใช้ในห้องปฏิบัติการ หรือ จองเพื่อซ่อมบำรุง \\
StartTime         & DATETIME &    & เวลาที่ผู้ยืมขอรับทรัพยากร \\
EndTime           & DATETIME &    & เวลาที่ผู้ยืมขอคืนทรัพยากร \\
ApproveStatus     & INT      & FK & สถานะการอนุมัติคำขอ อ้างอิง \texttt{ApproveStatus} \\
ApprovedBy        & INT      & FK & บัญชีผู้ใช้ที่เป็นผู้อนุมัติ \\
ReservationExpire & DATETIME &    & เวลาที่คำขอหมดอายุหากผู้จองไม่มารับทรัพยากรตามกำหนด \\
ActionTime        & DATETIME &    & เวลาที่ส่งคำขอ \\
ResolvedAt        & DATETIME &    & เวลาที่คำขอถูกอนุมัติ ปฏิเสธ หรือยกเลิก \\
\bottomrule
\end{tabularx}
\end{keeptable}

\begin{keeptable}
\noindent\textbf{ApproveStatus} — สถานะการอนุมัติ ใช้ร่วมกันทั้งคำขอจอง คำขอต่อเวลา และคำขออุทธรณ์

\begin{tabularx}{\linewidth}{@{} L{4.3cm} L{2.1cm} L{1.0cm} Y @{}}
\toprule
\hrow \thead{ฟิลด์} & \thead{ชนิดข้อมูล} & \thead{คีย์} & \thead{คำอธิบาย} \\
\midrule
ApproveStatusKey & INT     & PK & Surrogate Key ของตาราง \\
StatusName       & VARCHAR &    & ชื่อสถานะ ได้แก่ รออนุมัติ (Pending) / อนุมัติแล้ว (Approved) / ปฏิเสธ (Rejected) / ยกเลิก (Cancelled) \\
\bottomrule
\end{tabularx}
\end{keeptable}

\subsubsection*{กลุ่มการใช้งาน (Usage)}

\begin{keeptable}
\noindent\textbf{UsageLog} — บันทึกการส่งมอบและรับคืนทรัพยากร เป็นตารางหลักของรอบการยืมหนึ่งรอบ

\begin{tabularx}{\linewidth}{@{} L{4.3cm} L{2.1cm} L{1.0cm} Y @{}}
\toprule
\hrow \thead{ฟิลด์} & \thead{ชนิดข้อมูล} & \thead{คีย์} & \thead{คำอธิบาย} \\
\midrule
UsageKey          & INT      & PK & Surrogate Key ของตาราง \\
ReservationKey    & INT      & FK & การจองที่เป็นต้นทางของการใช้งานนี้ เว้นเป็น NULL ได้หากผู้ใช้ไม่ได้จองล่วงหน้า \\
AccountKey        & INT      & FK & บัญชีผู้ใช้ที่รับทรัพยากรไปใช้ \\
ResourceKey       & INT      & FK & ทรัพยากรที่ถูกนำไปใช้ \\
CurrentStatus     & INT      & FK & สถานะปัจจุบันของทรัพยากรในรอบการใช้งานนี้ \\
DueTime           & DATETIME &    & กำหนดเวลาที่ต้องคืน ปรับค่าใหม่ทุกครั้งที่การขอต่อเวลาได้รับอนุมัติ \\
PendingExtension  & INT      & FK & คำขอต่อเวลาที่ยังรอการอนุมัติ อ้างอิง \texttt{ExtensionRequest} \\
CheckoutTime      & DATETIME &    & เวลาที่ส่งมอบทรัพยากรให้ผู้ใช้ \\
CheckinTime       & DATETIME &    & เวลาที่รับทรัพยากรคืนจากผู้ใช้ \\
CheckoutCondition & INT      & FK & สภาพของทรัพยากรตอนส่งมอบให้ผู้ใช้ \\
CheckinCondition  & INT      & FK & สภาพของทรัพยากรตอนรับคืนจากผู้ใช้ \\
\bottomrule
\end{tabularx}
\end{keeptable}

\begin{keeptable}
\noindent\textbf{ExtensionRequest} — คำขอต่อเวลาการยืม

\begin{tabularx}{\linewidth}{@{} L{4.3cm} L{2.1cm} L{1.0cm} Y @{}}
\toprule
\hrow \thead{ฟิลด์} & \thead{ชนิดข้อมูล} & \thead{คีย์} & \thead{คำอธิบาย} \\
\midrule
ExtensionKey     & INT      & PK & Surrogate Key ของตาราง \\
UsageKey         & INT      & FK & รอบการใช้งานที่ขอต่อเวลา \\
RequestedBy      & INT      & FK & บัญชีผู้ใช้ที่ยื่นคำขอ \\
ExtendNo         & INT      &    & ลำดับครั้งที่ของการขอต่อเวลา ใช้ตรวจว่าถึงเพดาน \texttt{MaxExtendTime} แล้วหรือยัง \\
PreviousDueTime  & DATETIME &    & กำหนดคืนเดิมก่อนการต่อเวลา \\
RequestedDueTime & DATETIME &    & กำหนดคืนใหม่ที่ผู้ใช้ร้องขอ \\
ApproveStatus    & INT      & FK & สถานะการอนุมัติคำขอ \\
ApprovedBy       & INT      & FK & บัญชีผู้ใช้ที่เป็นผู้อนุมัติ \\
ActionTime       & DATETIME &    & เวลาที่ส่งคำขอ \\
ResolvedAt       & DATETIME &    & เวลาที่คำขอถูกอนุมัติ ปฏิเสธ หรือยกเลิก \\
\bottomrule
\end{tabularx}
\end{keeptable}

\subsubsection*{กลุ่มการตรวจสอบ (Inspection)}

\begin{keeptable}
\noindent\textbf{Inspection} — การตรวจสอบทรัพยากรโดยเจ้าหน้าที่หลังการใช้งานเสร็จสิ้น

\begin{tabularx}{\linewidth}{@{} L{4.3cm} L{2.1cm} L{1.0cm} Y @{}}
\toprule
\hrow \thead{ฟิลด์} & \thead{ชนิดข้อมูล} & \thead{คีย์} & \thead{คำอธิบาย} \\
\midrule
InspectionKey & INT      & PK & Surrogate Key ของตาราง \\
UsageKey      & INT      & FK & รอบการใช้งานที่ถูกตรวจสอบ \\
ResourceKey   & INT      & FK & ทรัพยากรที่ถูกตรวจสอบ \\
InspectorKey  & INT      & FK & บัญชีเจ้าหน้าที่ผู้ทำการตรวจสอบ \\
ConditionKey  & INT      & FK & ผลการตรวจสอบ อ้างอิง \texttt{ConditionLog} \\
AppealKey     & INT      & FK & การอุทธรณ์ที่เป็นต้นเหตุของการตรวจสอบรอบนี้ หากไม่ได้มาจากการอุทธรณ์ให้เป็น NULL \\
PenaltyKey    & INT      & FK & การหักคะแนนเครดิตที่เกิดจากการตรวจสอบนี้ \\
ActionTime    & DATETIME &    & เวลาที่ทำการตรวจสอบ \\
Notes         & TEXT     &    & บันทึกเพิ่มเติม \\
\bottomrule
\end{tabularx}
\end{keeptable}

\subsubsection*{กลุ่มบทลงโทษ (Penalty)}

\begin{keeptable}
\noindent\textbf{PenaltyInfo} — ข้อมูลการหักคะแนนเครดิตที่เกิดขึ้นจริงกับผู้ใช้แต่ละราย

\begin{tabularx}{\linewidth}{@{} L{4.3cm} L{2.1cm} L{1.0cm} Y @{}}
\toprule
\hrow \thead{ฟิลด์} & \thead{ชนิดข้อมูล} & \thead{คีย์} & \thead{คำอธิบาย} \\
\midrule
PenaltyKey     & INT      & PK & Surrogate Key ของตาราง \\
AccountKey     & INT      & FK & บัญชีผู้ใช้ที่ถูกหักคะแนน \\
UsageKey       & INT      & FK & รอบการใช้งานที่เป็นต้นเหตุของการหักคะแนน \\
Reason         & VARCHAR  &    & เหตุผลในการหักคะแนน \\
CreditDeducted & INT      &    & จำนวนคะแนนเครดิตที่ถูกหัก \\
ActionTime     & DATETIME &    & เวลาที่ถูกหักคะแนน \\
ExpirationTime & DATETIME &    & วันที่การหักคะแนนหมดอายุ \\
Appealed       & BOOL     &    & การอุทธรณ์ได้รับการอนุมัติแล้วหรือไม่ \\
InEffect       & BOOL     &    & บทลงโทษยังมีผลบังคับอยู่หรือไม่ โดยหลักใช้คู่กับ \texttt{ExpirationTime} \\
\bottomrule
\end{tabularx}
\end{keeptable}

\subsubsection*{กลุ่มการอุทธรณ์ (Appeal)}

\begin{keeptable}
\noindent\textbf{AppealInfo} — ข้อมูลการขออุทธรณ์ผลการประเมินความเสียหาย

\begin{tabularx}{\linewidth}{@{} L{4.3cm} L{2.1cm} L{1.0cm} Y @{}}
\toprule
\hrow \thead{ฟิลด์} & \thead{ชนิดข้อมูล} & \thead{คีย์} & \thead{คำอธิบาย} \\
\midrule
AppealKey       & INT      & PK & Surrogate Key ของตาราง \\
OriginalPenalty & INT      & FK & การหักคะแนนเดิมที่ถูกอุทธรณ์ \\
NewPenalty      & INT      & FK & การหักคะแนนใหม่ ในกรณีที่ผลอุทธรณ์ทำให้บทลงโทษลดลง \\
FiledBy         & INT      & FK & บัญชีผู้ใช้ที่ยื่นอุทธรณ์ \\
ResolvedBy      & INT      & FK & บัญชีผู้ใช้ที่เป็นผู้ตรวจสอบใหม่ \\
AppealReason    & TEXT     &    & เหตุผลในการขออุทธรณ์ \\
ApproveStatus   & INT      & FK & สถานะการพิจารณาคำขอ \\
ActionTime      & DATETIME &    & เวลาที่ยื่นอุทธรณ์ \\
ResolvedAt      & DATETIME &    & เวลาที่คำขอถูกอนุมัติ ปฏิเสธ หรือยกเลิก \\
\bottomrule
\end{tabularx}
\end{keeptable}

\subsubsection*{อื่น ๆ (Other)}

\begin{keeptable}
\noindent\textbf{Image} — ข้อมูลรูปภาพที่ถ่ายไว้ประกอบการตรวจสอบทรัพยากร

\begin{tabularx}{\linewidth}{@{} L{4.3cm} L{2.1cm} L{1.0cm} Y @{}}
\toprule
\hrow \thead{ฟิลด์} & \thead{ชนิดข้อมูล} & \thead{คีย์} & \thead{คำอธิบาย} \\
\midrule
ImageKey          & INT     & PK & Surrogate Key ของตาราง \\
SubmittedBy       & INT     & FK & บัญชีผู้ใช้ที่เป็นผู้อัปโหลดภาพ \\
UsageKey          & INT     & FK & รอบการใช้งานที่ภาพนี้สังกัด \\
ResourceKey       & INT     & FK & ทรัพยากรที่ปรากฏในภาพ \\
InspectionKey     & INT     & FK & การตรวจสอบที่ภาพนี้สังกัด หากเป็นภาพที่ผู้ยืมถ่ายเองให้เป็น NULL \\
ImageURL          & VARCHAR &    & ชื่อไฟล์หรือลิงก์ของไฟล์ภาพ \\
SubmissionTypeKey & INT     & FK & ประเภทของการส่งภาพ อ้างอิง \texttt{ImageSubmissionType} \\
\bottomrule
\end{tabularx}
\end{keeptable}

\begin{keeptable}
\noindent\textbf{ImageSubmissionType} — ประเภทของการส่งภาพ

\begin{tabularx}{\linewidth}{@{} L{4.3cm} L{2.1cm} L{1.0cm} Y @{}}
\toprule
\hrow \thead{ฟิลด์} & \thead{ชนิดข้อมูล} & \thead{คีย์} & \thead{คำอธิบาย} \\
\midrule
SubmissionTypeKey & INT     & PK & Surrogate Key ของตาราง \\
SubmissionType    & VARCHAR &    & ชื่อประเภท ได้แก่ Inspection (ภาพจากการตรวจสอบโดยเจ้าหน้าที่) / BeforeBorrow (ภาพก่อนยืม) / AfterBorrow (ภาพหลังยืม) \\
\bottomrule
\end{tabularx}
\end{keeptable}

\begin{keeptable}
\noindent\textbf{Notification} — ข้อมูลการแจ้งเตือนที่ส่งถึงผู้ใช้

\begin{tabularx}{\linewidth}{@{} L{4.3cm} L{2.1cm} L{1.0cm} Y @{}}
\toprule
\hrow \thead{ฟิลด์} & \thead{ชนิดข้อมูล} & \thead{คีย์} & \thead{คำอธิบาย} \\
\midrule
NotificationKey     & INT      & PK & Surrogate Key ของตาราง \\
AccountKey          & INT      & FK & บัญชีผู้ใช้ที่ได้รับข้อความ \\
NotificationType    & INT      & FK & ประเภทของการแจ้งเตือน \\
NotificationContent & TEXT     &    & เนื้อความที่ส่ง \\
SentTime            & DATETIME &    & เวลาที่ส่งข้อความ \\
IsRead              & BOOL     &    & ผู้ใช้เปิดอ่านข้อความแล้วหรือยัง \\
\bottomrule
\end{tabularx}
\end{keeptable}

\begin{keeptable}
\noindent\textbf{NotificationType} — ประเภทของการแจ้งเตือน

\begin{tabularx}{\linewidth}{@{} L{4.3cm} L{2.1cm} L{1.0cm} Y @{}}
\toprule
\hrow \thead{ฟิลด์} & \thead{ชนิดข้อมูล} & \thead{คีย์} & \thead{คำอธิบาย} \\
\midrule
NotificationTypeKey  & INT     & PK & Surrogate Key ของตาราง \\
NotificationTypeName & VARCHAR &    & ชื่อประเภท เช่น เตือนใกล้ครบกำหนดคืน / แจ้งผลการอนุมัติ \\
\bottomrule
\end{tabularx}
\end{keeptable}

\begin{keeptable}
\noindent\textbf{RepairLog} — ข้อมูลการซ่อมและระยะเวลาที่ทรัพยากรไม่พร้อมให้บริการ

\begin{tabularx}{\linewidth}{@{} L{4.3cm} L{2.1cm} L{1.0cm} Y @{}}
\toprule
\hrow \thead{ฟิลด์} & \thead{ชนิดข้อมูล} & \thead{คีย์} & \thead{คำอธิบาย} \\
\midrule
RepairKey             & INT      & PK & Surrogate Key ของตาราง \\
ReservationKey        & INT      & FK & การจองช่วงเวลาซ่อมล่วงหน้า หากมีการจองไว้ \\
RepairedBy            & INT      & FK & บัญชีผู้ใช้ที่เป็นผู้ดำเนินการซ่อม \\
ResourceKey           & INT      & FK & ทรัพยากรที่ถูกซ่อม \\
BeginRepairDate       & DATETIME &    & เวลาที่เริ่มการซ่อมตามจริง \\
EndRepairDate         & DATETIME &    & เวลาที่การซ่อมเสร็จสิ้นตามจริง \\
ConditionBeforeRepair & INT      & FK & สภาพของทรัพยากรก่อนการซ่อม \\
ConditionAfterRepair  & INT      & FK & สภาพของทรัพยากรหลังการซ่อม \\
\bottomrule
\end{tabularx}
\end{keeptable}

\subsubsection*{การเปลี่ยนแปลงจากการออกแบบรอบก่อนหน้า}

การออกแบบชุดนี้แทนที่การออกแบบรอบก่อนหน้าทั้งหมด ตารางเดิมเจ็ดตารางถูกยกเลิก และหน้าที่ของแต่ละตารางถูกกระจายไปยังตารางใหม่ดังนี้

\begin{keeptable}
\noindent\textbf{ตารางเทียบโครงสร้างเดิมกับโครงสร้างใหม่}

\begin{tabularx}{\linewidth}{@{} L{4.3cm} Y @{}}
\toprule
\hrow \thead{ตารางเดิม} & \thead{หน้าที่ในการออกแบบรอบนี้} \\
\midrule
StockInfo        & รวมเข้ากับ \texttt{ResourceInfo} ซึ่งเป็นตารางกลางของทั้งอุปกรณ์และห้อง \\
BorrowInfo       & แยกเป็น \texttt{Reservations} (คำขอจอง) และ \texttt{UsageLog} (การส่งมอบและรับคืนจริง) \\
RequestLog       & ยกเลิก การยืมหลายรายการในคำขอเดียวใช้หลายระเบียนใน \texttt{Reservations} แทน \\
ExtendBorrow     & เปลี่ยนเป็น \texttt{ExtensionRequest} ซึ่งเก็บกำหนดคืนเดิมและกำหนดคืนใหม่แยกกัน \\
UserExamination  & รวมเป็น \texttt{ConditionLog} ร่วมกับการตรวจสอบของเจ้าหน้าที่ \\
StaffExamination & แยกเป็น \texttt{Inspection} (ตัวการตรวจสอบ) และ \texttt{ConditionLog} (ผลสภาพที่บันทึก) \\
FacultyPermission & เปลี่ยนเป็น \texttt{Authority} ร่วมกับ \texttt{AuthorityRole} และ \texttt{ManagementGroup} \\
\bottomrule
\end{tabularx}
\end{keeptable}

\subsubsection*{หมายเหตุการออกแบบ}

\begin{itemize}
  \item \textbf{ประเด็นที่ต้องตกลงร่วมกับส่วนเทคโนโลยีที่ใช้} การออกแบบรอบนี้ไม่มีฟิลด์ชนิดรายการ (LIST) เหลืออยู่เลย ตารางที่เคยเป็นเหตุผลของการเลือกใช้ MongoDB ได้แก่ \texttt{FacultyPermission}, \texttt{BorrowRule} และ \texttt{Image} ถูกแตกเป็นตารางย่อยแบบ Normalize ทั้งหมดแล้ว หัวข้อสถาปัตยกรรมระบบและเทคโนโลยีที่ใช้จึงยังคงระบุ MongoDB ไว้ตามเอกสารนำเสนอเดิม แต่ยังไม่มีข้อมูลส่วนใดในการออกแบบรอบนี้ที่ต้องจัดเก็บบน MongoDB ควรทบทวนประเด็นนี้ร่วมกันในรอบถัดไป
  \item ตาราง \texttt{ResourceInfo} และ \texttt{ConditionLog} อ้างอิงถึงกันทั้งสองทาง (\texttt{ResourceInfo.ConditionKey} ชี้ไปยัง \texttt{ConditionLog} ขณะที่ \texttt{ConditionLog.ResourceKey} ชี้กลับมายัง \texttt{ResourceInfo}) ในทางปฏิบัติจึงต้องกำหนดให้ฝั่งใดฝั่งหนึ่งรับค่า NULL ได้ในตอนสร้างระเบียนแรก หรือใช้ข้อบังคับแบบเลื่อนการตรวจสอบ (Deferred Constraint)
  \item ชื่อฟิลด์ \texttt{RoleKey} ถูกใช้กับสองความหมาย โดยใน \texttt{AccountInfo} หมายถึงบทบาทระดับระบบ (อ้างอิง \texttt{RoleInfo}) ส่วนใน \texttt{Eligibility} หมายถึงตำแหน่งภายในกลุ่ม (อ้างอิง \texttt{AuthorityRole}) ควรตั้งชื่อให้ต่างกันเพื่อลดความสับสนในการพัฒนา
  \item การตั้งชื่อฟิลด์ที่เป็น Foreign Key ยังไม่สม่ำเสมอ บางฟิลด์ลงท้ายด้วย \texttt{Key} เช่น \texttt{ManagementGroupKey} ขณะที่บางฟิลด์ไม่ลงท้าย เช่น \texttt{FacultyInfo.ManagementGroup}, \texttt{ResourceInfo.BorrowRule} และ \texttt{PenaltyRule.PenaltyType} ควรปรับให้เป็นรูปแบบเดียวกันก่อนลงมือสร้างฐานข้อมูลจริง
  \item การจัดกลุ่มในพจนานุกรมข้อมูลอ้างตามหัวข้อในเอกสารต้นทาง ยกเว้นกลุ่ม ``ทรัพยากรและสภาพ'' และ ``กฎการยืมและบทลงโทษ'' ซึ่งเอกสารต้นทางไม่ได้กำกับหัวข้อไว้ รายงานนี้จึงเพิ่มขึ้นเพื่อให้อ่านได้ต่อเนื่อง
  \item ข้อคลาดเคลื่อนเล็กน้อยในเอกสารต้นทางที่ปรับให้ถูกต้องแล้วในตารางข้างต้น ได้แก่ \texttt{UsageLog.AcocuntKey} แก้เป็น \texttt{AccountKey}, \texttt{UsageLog.CheckinTIme} แก้เป็น \texttt{CheckinTime}, \texttt{ClubInfo.ClubName} ระบุชนิดเป็น INT แก้เป็น VARCHAR เนื่องจากเก็บชื่อชมรม และ \texttt{Inspection.AppealKey} ระบุปลายทางเป็น \texttt{Appeal} แก้เป็น \texttt{AppealInfo} ตามชื่อตารางจริง ส่วนหัวข้อ ``Images'' ในเอกสารต้นทางใช้ชื่อตามแผนภาพ ERD ว่า \texttt{Image}
\end{itemize}

\clearpage
\subsection{รูปแบบการเรียกใช้งานหลังบ้านและงานเบื้องหลัง (Backend Interaction \& Background Jobs)}

ระบบออกแบบรูปแบบการรับ–ส่งข้อมูลระหว่าง Client กับ Backend ออกเป็นสามกลุ่มตามลักษณะการใช้งาน เพื่อให้ข้อมูลที่ต้องแข่งขันแบบเรียลไทม์เป็นปัจจุบันอยู่เสมอ ขณะเดียวกันก็ลดภาระการเรียกข้อมูลที่ไม่จำเป็นสำหรับข้อมูลที่เปลี่ยนแปลงไม่บ่อย

\subsubsection*{กลุ่มที่ 1 — ข้อมูลเรียลไทม์ (Polling)}

ข้อมูลที่เปลี่ยนแปลงบ่อยและต้องแสดงผลใกล้เคียงเรียลไทม์ ระบบใช้การดึงข้อมูลซ้ำตามรอบเวลา (Polling ด้วย TanStack Query) เนื่องจากบางรายการมีการแข่งขันกันระหว่างผู้ใช้ เช่น การกดยืมหรือจองตัดหน้ากัน

\begin{tabularx}{\dimexpr\textwidth-\leftskip\relax}{@{} Y L{4.6cm} L{2.3cm} @{}}
\toprule
\hrow \thead{ข้อมูลที่ดึงซ้ำ} & \thead{เหตุผล} & \thead{รอบเวลา} \\
\midrule
ของคงเหลือ / วันที่พร้อมให้ยืมของอุปกรณ์ที่กำลังดู & มีโอกาสถูกผู้อื่นกดยืมตัดหน้า & 10–15 วินาที \\
สล็อตเวลาว่างของห้อง/สถานที่ (T3) & การจองมักแข่งขันกันแบบเรียลไทม์ & 10–15 วินาที \\
สถานะคำขอยืมของตนเองที่ยัง ``รออนุมัติ'' & รอผลอนุมัติจากระบบหรืออาจารย์ & 30 วินาที \\
การแจ้งเตือนในแอป (Notification) & แจ้งการอนุมัติ เตือนคืน และผลอุทธรณ์ทันที & 30–60 วินาที \\
คิวงานฝั่งเจ้าหน้าที่ (คำขอรอจัดเตรียม, ของรอรับคืน) & จัดการคำขอใหม่ได้อย่างรวดเร็ว & 30 วินาที \\
คิวรออนุมัติฝั่งอาจารย์ (T2, เครดิตต่ำ, ต่ออายุ, อุทธรณ์) & ตอบสนองคำขอใหม่ได้ทันท่วงที & 30–60 วินาที \\
\bottomrule
\end{tabularx}

\subsubsection*{กลุ่มที่ 2 — เรียกเมื่อใช้งาน (Fetch \& Refetch)}

ข้อมูลที่เรียกเพียงครั้งเดียวเมื่อเปิดใช้งาน และเรียกซ้ำ (refetch) เฉพาะหลังการกระทำที่เกี่ยวข้อง แบ่งเป็นฝั่งอ่านข้อมูล (Read) และฝั่งการกระทำ (Action / Mutation)

\noindent\textbf{ฝั่งอ่านข้อมูล (Read)}
\begin{enumerate}
  \item รายการอุปกรณ์ (ค้นหา + กรอง + แบ่งหน้า)
  \item รายละเอียดอุปกรณ์รายชิ้น
  \item โปรไฟล์ตัวเอง + บทบาท (Role) + สังกัด
  \item เครดิตตัวเอง + ระดับเครดิต + รายการหักที่ยัง active
  \item รายการที่กำลังยืมอยู่ + ประวัติการยืม
  \item รายละเอียดการยืม (Loan): รูปตอนรับ และสิทธิ์ต่ออายุออนไลน์คงเหลือ
  \item สถานะและรายละเอียดการอุทธรณ์
  \item Dashboard / รายงานสถิติ (สำหรับเจ้าหน้าที่และผู้ดูแลระบบ)
  \item Master data: หมวดหมู่ ระดับอุปกรณ์ (Tier) และระดับความเสียหาย — เก็บ cache ระยะยาว (staleTime ยาว)
\end{enumerate}

\noindent\textbf{ฝั่งการกระทำ (Action / Mutation)} — เมื่อทำรายการเสร็จ ระบบจะเรียกข้อมูลที่เกี่ยวข้องซ้ำ (refetch)
\begin{itemize}
  \item ผู้ยืม: ส่งคำขอยืม, ยกเลิกคำขอ, ยืนยันรับของ + อัปโหลดรูปตอนรับ, แจ้งคืน + อัปโหลดรูปตอนคืน, ขอต่ออายุ, ยื่นอุทธรณ์
  \item เจ้าหน้าที่: จัดเตรียม/เลือกชิ้น, บันทึกรับ–คืน, ประเมินความเสียหาย, จัดการอุปกรณ์
  \item อาจารย์: อนุมัติ/ปฏิเสธคำขอ, ตัดสินอุทธรณ์
  \item ผู้ดูแลระบบ: จัดการผู้ใช้/สิทธิ์, ตั้งค่าระบบ
\end{itemize}

\subsubsection*{กลุ่มที่ 3 — งานเบื้องหลังตามเวลา (Scheduled / Cron Jobs)}

งานที่ระบบทำงานอัตโนมัติตามรอบเวลาที่ฝั่ง Backend โดยส่วนใหญ่เป็นงานรอบรายวัน

\begin{enumerate}
  \item \textbf{Mark overdue:} เลยกำหนดคืน $\rightarrow$ เปลี่ยนสถานะและหักเครดิตรายวัน
  \item \textbf{Mark lost:} เกินกำหนด 2 สัปดาห์ $\rightarrow$ ถือว่าหายและหักเครดิตเพิ่ม
  \item \textbf{หมดอายุบทลงโทษ:} ครบอายุโทษ $\rightarrow$ คืนคะแนนและอัปเดตระดับเครดิต
  \item \textbf{เตือนใกล้ครบกำหนด:} สร้าง notification และส่ง email ล่วงหน้า 1–2 วัน
  \item \textbf{คำนวณวันที่พร้อมให้ยืม:} อัปเดตวันพร้อมให้ยืมของแต่ละประเภท/ชิ้น
  \item \textbf{สร้างรอบตรวจสถานที่รายวัน (T3):} เปิด task ให้เจ้าหน้าที่ตรวจห้อง
  \item \textbf{สรุปสถิติรายวัน:} อัปเดตข้อมูล Dashboard และรายงาน
\end{enumerate}

\noindent นอกจากนี้ยังมีงานรอบรายชั่วโมง คือการยกเลิกคำขอที่หมดอายุการรับของ (คำขอที่ไม่มีผู้มารับภายในกำหนดจะถูกยกเลิกอัตโนมัติทุกชั่วโมง)
